\documentclass[11pt,a4paper,oneside]{article}

% Standalone ProveIt research note.  No external TeX inputs are required.
\usepackage[a4paper,margin=25mm,headheight=15pt]{geometry}
\usepackage[T1]{fontenc}
\usepackage[utf8]{inputenc}
\IfFileExists{libertinus.sty}{\usepackage{libertinus}}{\usepackage{lmodern}}
\usepackage{microtype}
\usepackage{amsmath,amssymb,amsthm,mathtools,mathrsfs,bm}
\usepackage{booktabs,longtable,tabularx,array}
\usepackage{graphicx}
\usepackage{xcolor}
\usepackage{enumitem}
\usepackage{fancyhdr}
\usepackage{xurl}
\usepackage{listings}
\usepackage[most]{tcolorbox}
\usepackage{hyperref}
\usepackage[nameinlink,noabbrev,capitalise]{cleveref}

\definecolor{linkblue}{RGB}{31,76,122}
\definecolor{deepolive}{RGB}{67,79,45}
\definecolor{softolive}{RGB}{236,239,226}
\definecolor{darkteal}{RGB}{31,83,84}
\definecolor{softteal}{RGB}{232,244,243}
\definecolor{warmgray}{RGB}{92,87,78}
\definecolor{softcoral}{RGB}{246,225,219}
\definecolor{softgold}{RGB}{251,245,230}
\definecolor{rulegray}{RGB}{195,201,208}
\definecolor{shadegray}{RGB}{246,247,249}

\hypersetup{
  hypertexnames=false,
  colorlinks=true,
  linkcolor=deepolive,
  citecolor=darkteal,
  urlcolor=linkblue,
  pdftitle={The Double Factorial and Its Inverse: Complete Asymptotic Transseries},
  pdfauthor={Vladimir Reshetnikov},
  pdfsubject={Parity-resolved double factorial asymptotics, Borel summation, Lambert-W inverse transseries, Bell-polynomial coefficient formulae, and the periodic OEIS interpolation},
  pdfkeywords={double factorial, inverse double factorial, asymptotic transseries, Lambert W, Bernoulli numbers, Bell polynomials, Lagrange inversion, Borel summation}
}

\pagestyle{fancy}
\fancyhf{}
\fancyhead[L]{\small Double factorial transseries}
\fancyhead[R]{\small\nouppercase{\leftmark}}
\fancyfoot[C]{\thepage}
\renewcommand{\headrulewidth}{0.4pt}

\setcounter{tocdepth}{2}
\setcounter{secnumdepth}{3}
\setlist[itemize]{leftmargin=1.5em,itemsep=0.25ex,topsep=0.6ex}
\setlist[enumerate]{leftmargin=1.9em,itemsep=0.35ex,topsep=0.6ex}
\allowdisplaybreaks[2]
\emergencystretch=2em

\theoremstyle{plain}
\newtheorem{theorem}{Theorem}[section]
\newtheorem{proposition}[theorem]{Proposition}
\newtheorem{lemma}[theorem]{Lemma}
\newtheorem{corollary}[theorem]{Corollary}
\theoremstyle{definition}
\newtheorem{definition}[theorem]{Definition}
\newtheorem{algorithm}[theorem]{Algorithm}
\newtheorem{example}[theorem]{Example}
\theoremstyle{remark}
\newtheorem{remark}[theorem]{Remark}

\crefname{algorithm}{algorithm}{algorithms}
\Crefname{algorithm}{Algorithm}{Algorithms}

\newtcolorbox{mainbox}[1][]{%
  enhanced,breakable,colback=softolive,colframe=deepolive,
  boxrule=0.7pt,arc=1.2mm,left=2mm,right=2mm,top=1.4mm,bottom=1.4mm,
  title={Main result},fonttitle=\bfseries,#1}
\newtcolorbox{methodbox}[1][]{%
  enhanced,breakable,colback=softteal,colframe=darkteal,
  boxrule=0.6pt,arc=1.2mm,left=2mm,right=2mm,top=1.2mm,bottom=1.2mm,
  title={Coefficient calculus},fonttitle=\bfseries,#1}
\newtcolorbox{warningbox}[1][]{%
  enhanced,breakable,colback=softcoral,colframe=warmgray,
  boxrule=0.6pt,arc=1.2mm,left=2mm,right=2mm,top=1.2mm,bottom=1.2mm,
  title={Interpretive point},fonttitle=\bfseries,#1}
\newtcolorbox{summarybox}[1][]{%
  enhanced,breakable,colback=softgold,colframe=warmgray,
  boxrule=0.6pt,arc=1.2mm,left=2mm,right=2mm,top=1.2mm,bottom=1.2mm,
  title={Summary},fonttitle=\bfseries,#1}

\lstdefinestyle{pseudo}{%
  basicstyle=\ttfamily\small,
  columns=fullflexible,
  frame=single,
  rulecolor=\color{rulegray},
  backgroundcolor=\color{shadegray},
  showstringspaces=false,
  keepspaces=true,
  breaklines=true,
  xleftmargin=0.4em,xrightmargin=0.4em,
  aboveskip=0.8em,belowskip=0.8em
}
\lstset{style=pseudo}

\newcommand{\e}{\mathrm e}
\newcommand{\dd}{\mathop{}\!\mathrm d}
\newcommand{\ii}{\mathrm i}
\newcommand{\R}{\mathbb R}
\newcommand{\C}{\mathbb C}
\newcommand{\Q}{\mathbb Q}
\newcommand{\N}{\mathbb N}
\newcommand{\W}{\operatorname W}
\newcommand{\Log}{\operatorname{Log}}
\newcommand{\coeff}[2]{\left[#1^{#2}\right]}
\newcommand{\BellP}[2]{\mathrm B_{#1,#2}}
\newcommand{\BellC}[1]{\mathrm Y_{#1}}
\newcommand{\stirlingone}[2]{\genfrac{[}{]}{0pt}{}{#1}{#2}}
\newcommand{\D}{\mathcal D}
\newcommand{\DO}{\mathfrak D}
\newcommand{\etaD}{\eta}
\newcommand{\asympseq}{\mathrel{\sim}}
\DeclareMathOperator{\Res}{Res}
\DeclareMathOperator{\sgn}{sgn}

\begin{document}

\title{\textbf{The Double Factorial and Its Inverse}\\[0.35em]
  \Large Complete Asymptotic Transseries, Borel Summation,\\
  and General Coefficient Formulae}
\author{Vladimir Reshetnikov\\
  \small ProveIt research note}
\date{3 September 2026}
\maketitle

\begin{abstract}
The double factorial sequence $n!!$ is not the restriction of one canonical
slowly varying function: its even and odd subsequences are two gamma-function
branches whose quotient contains a non-vanishing parity factor.  This article
therefore treats three related inverse problems separately and exactly: the two
parity-resolved smooth inverses, the integer-valued generalized inverse, and the
eventual inverse of the periodic analytic continuation recorded in OEIS
A006882.

For $r\in\{0,1\}$ we introduce
\[
 \D_r(x)=\kappa_r2^{x/2}\Gamma\!\left(\frac{x}{2}+1\right),
 \qquad \kappa_0=1,\qquad \kappa_1=\sqrt{\frac2\pi},
\]
so that $\D_r(n)=n!!$ whenever $n\equiv r\pmod2$.  With $s=x+1$ and
$C_0=\tfrac12\log\pi$, $C_1=\tfrac12\log2$, the optimally centred logarithm is
\[
 \log\D_r(s-1)\sim \frac{s}{2}(\log s-1)+C_r
 +\sum_{k\ge1}\frac{\beta_k}{s^{2k-1}},
 \qquad
 \beta_k=\frac{(1-2^{2k-1})B_{2k}}{2k(2k-1)}.
\]
We give all-order additive and multiplicative coefficient formulae in terms of
Bernoulli numbers, integer partitions, and complete Bell polynomials.  More
strongly, the divergent logarithmic series is the asymptotic expansion of the
exact Borel--Laplace integral
\[
 \int_0^\infty \e^{-st}
 \left(\frac{1}{2t\sinh t}-\frac{1}{2t^2}\right)\dd t.
\]
This supplies alternating remainder bounds, the complete Borel singularity
lattice, factorial large-order estimates, and the optimal-truncation scale.

For the large real inverse, put
\[
 R_r=\log y-C_r,\qquad w_r=\W_0(2R_r/\e),\qquad
 T_r=\frac{2R_r}{w_r},\qquad \Lambda_r=w_r+1=\log T_r.
\]
Then, for every fixed $N$,
\[
 \D_r^{-1}(y)=T_r-1+
 \sum_{n=1}^{N}\frac{d_{2n-1}(\Lambda_r)}{T_r^{2n-1}}
 +O\!\left(\frac{1}{T_r^{2N+1}\Lambda_r}\right).
\]
The coefficient $d_{2n-1}$ is a rational polynomial in $\Lambda^{-1}$ of
degree at most $2n-1$.  We derive both a triangular recurrence and a finite,
non-recursive Lagrange--B\"urmann formula.  The latter is expressed through
partial Bell polynomials and proves the global sign law
$(-1)^{n+1}d_{2n-1}\in\Lambda^{-1}\Q_{>0}[\Lambda^{-1}]$.  Explicit blocks are
listed through $T^{-9}$, and the Lambert core is expanded to all orders in
iterated logarithms using unsigned Stirling numbers of the first kind.

Finally, for the OEIS continuation
\[
 \DO(x)=2^{x/2}\Gamma\!\left(\frac{x}{2}+1\right)
 \left(\frac2\pi\right)^{(1-\cos\pi x)/4},
\]
the inverse has a leading log-periodic sector of order $1/\log T$, larger than
all inverse-power corrections.  We give its closed Fourier--Lagrange
coefficients and a second triangular recurrence that couples this periodic
sector to the Bernoulli grid.  Thus the ambiguity in ``the inverse double
factorial'' is resolved rather than hidden.
\end{abstract}

\noindent\textbf{Keywords.}
Double factorial; inverse double factorial; Lambert $W$; asymptotic
transseries; Stirling series; Borel summation; Bernoulli numbers; Bell
polynomials; Lagrange--B\"urmann inversion; parity; log-periodic asymptotics.

{\small\tableofcontents}
\clearpage

\section{Introduction}
\label{sec:intro}

The double factorial is defined on nonnegative integers by
\begin{equation}
 0!!=1!!=1,\qquad n!!=n(n-2)!!\quad(n\ge2).
 \label{eq:sequence-def}
\end{equation}
Equivalently, $n!!$ is the product of the positive integers not exceeding $n$
that have the same parity as $n$.  The sequence begins
\[
 1,1,2,3,8,15,48,105,384,945,\ldots
\]
and is OEIS A006882 \cite{OEIS}.  Its elementary closed forms are
\begin{equation}
 (2m)!!=2^m m!,
 \qquad
 (2m-1)!!=\frac{(2m)!}{2^m m!}
          =\frac{2^m\Gamma(m+1/2)}{\sqrt\pi}.
 \label{eq:parity-closed}
\end{equation}
These identities already expose the main asymptotic issue: the even and odd
subsequences have different multiplicative constants of leading order.  A
single ordinary Poincar\'e expansion with constant coefficients cannot represent
both without retaining a parity monomial such as $(-1)^n$.

The word \emph{inverse} introduces a second issue.  A sequence has no ordinary
inverse function until an interpolation or a generalized-inverse convention is
chosen.  There are three natural choices here.
\begin{enumerate}
\item The even and odd subsequences have canonical gamma continuations
      $\D_0$ and $\D_1$.  Their smooth inverses are the cleanest objects for
      high-accuracy asymptotics and preserve the exact parity normalization.
\item The integer sequence has a right-continuous generalized inverse, obtained
      by applying a parity-preserving floor to the two smooth inverses.
\item OEIS records a single analytic continuation $\DO$ which agrees with both
      parities.  It contains a periodic factor, and its eventual inverse has a
      log-periodic transseries qualitatively different from either $\D_r^{-1}$.
\end{enumerate}
The article develops all three and makes their relationship explicit.

Two asymptotic scales must also be separated.  Taking the logarithm converts
the main growth into
\[
 \frac{s}{2}(\log s-1),\qquad s=x+1.
\]
Its inverse is exactly a Lambert-$W$ expression.  Expanding $W$ produces an
\emph{outer} series in powers of $1/\log\log y$ with polynomial dependence on
$\log\log\log y$.  The Bernoulli corrections generate a much finer grid in
odd inverse powers of the Lambert core $T$.  The fixed shift $-1$ lies between
these levels: it is beyond every finite order of the outer Lambert expansion,
but much larger than $1/(T\log T)$.  Treating all these terms as one ordinary
series loses the natural well-ordering of the transseries.

\begin{warningbox}
There is no contradiction between the two statements
\[
 n!!\sim A_n\sqrt n\,(n/\e)^{n/2}
 \quad\hbox{and}\quad
 n!!\sim A_n\bigl((n+1)/\e\bigr)^{(n+1)/2}.
\]
They use different normal forms.  The first leaves a factor $\sqrt n$ and has
an uncentred Stirling series.  The second absorbs that factor by the optimal
shift $s=n+1$ and has an odd inverse-power logarithmic series.  The centred
form is far better for inversion.
\end{warningbox}

\subsection{Principal formulae}

Put
\begin{equation}
 \kappa_0=1,\qquad \kappa_1=\sqrt{2/\pi},
 \qquad
 C_0=\frac12\log\pi,\qquad C_1=\frac12\log2.
 \label{eq:constants}
\end{equation}
For $r\in\{0,1\}$ define
\begin{equation}
 \D_r(x)=\kappa_r2^{x/2}\Gamma\!\left(\frac{x}{2}+1\right).
 \label{eq:Dr-def}
\end{equation}
The direct logarithmic coefficients in the centred variable $s=x+1$ are
\begin{equation}
 \beta_k=\frac{(1-2^{2k-1})B_{2k}}{2k(2k-1)}
 =(-1)^k\frac{(2k-2)!\,\etaD(2k)}{\pi^{2k}},
 \label{eq:beta-def}
\end{equation}
where $\etaD(z)=(1-2^{1-z})\zeta(z)$ is the Dirichlet eta function.  Let
\begin{equation}
 \alpha_k=2\beta_k
 =\frac{(1-2^{2k-1})B_{2k}}{k(2k-1)}.
 \label{eq:alpha-def}
\end{equation}

\begin{mainbox}
For $y\to+\infty$ and $r\in\{0,1\}$, define
\begin{equation}
 R_r=\log y-C_r,
 \qquad w_r=\W_0(2R_r/\e),
 \qquad T_r=\frac{2R_r}{w_r}=\e^{w_r+1},
 \qquad \Lambda_r=w_r+1.
 \label{eq:inverse-core-summary}
\end{equation}
Then
\begin{equation}
 \D_r^{-1}(y)
 \sim T_r-1+\sum_{n\ge1}
 \frac{d_{2n-1}(\Lambda_r)}{T_r^{2n-1}},
 \label{eq:inverse-main-summary}
\end{equation}
where $d_{2n-1}(\Lambda)\in
\Lambda^{-1}\Q[\Lambda^{-1}]$ and
\begin{align}
 d_1(\Lambda)&=\frac{1}{6\Lambda},
 \label{eq:d1-summary}\\
 d_3(\Lambda)&=-\frac{14\Lambda^2+10\Lambda+5}
 {360\Lambda^3},
 \label{eq:d3-summary}\\
 d_5(\Lambda)&=\frac{2232\Lambda^4+1176\Lambda^3+714\Lambda^2
 +350\Lambda+105}{45360\Lambda^5}.
 \label{eq:d5-summary}
\end{align}
Every coefficient is generated by the recurrence in
\cref{thm:branch-inverse}; the finite closed formula is
\cref{thm:closed-coefficients}.  Expanding the retained Lambert function gives
an outer iterated-logarithmic transseries with coefficients generated from
unsigned Stirling numbers; see \cref{sec:outer}.
\end{mainbox}

\section{Exact continuations and the meaning of the inverse}
\label{sec:continuations}

\subsection{The two parity-resolved gamma branches}

\begin{proposition}[Exact parity interpolation]
\label{prop:parity-interpolation}
For every integer $n\ge0$ and $r\in\{0,1\}$ with $n\equiv r\pmod2$,
\begin{equation}
 \D_r(n)=n!!.
 \label{eq:Dr-interpolation}
\end{equation}
Moreover, $\D_r$ is strictly increasing on $[0,\infty)$, and hence has a
single-valued inverse on its range.
\end{proposition}

\begin{proof}
If $n=2m$, then
$\D_0(2m)=2^m\Gamma(m+1)=2^m m!=(2m)!!$.  If $n=2m+1$, then
\[
 \D_1(2m+1)=\sqrt{\frac2\pi}\,2^{m+1/2}\Gamma(m+3/2)
 =\frac{2^{m+1}}{\sqrt\pi}\Gamma(m+3/2)=(2m+1)!!.
\]
For real $x\ge0$,
\begin{equation}
 \frac{\D_r'(x)}{\D_r(x)}
 =\frac12\left(\log2+\psi\!\left(\frac{x}{2}+1\right)\right),
 \label{eq:Dr-logder-exact}
\end{equation}
where $\psi=\Gamma'/\Gamma$.  The trigamma series
$\psi'(z)=\sum_{m\ge0}(z+m)^{-2}$ is positive for $z>0$, so the right-hand
side is increasing.  At $x=0$ it equals
$\tfrac12(\log2-\gamma)>0$.  Thus it remains positive.
\end{proof}

\begin{remark}
The factor $\kappa_r$ affects only the target normalization; it disappears
from the logarithmic derivative.  Consequently both parity branches have the
same local conditioning and the same inverse correction polynomials.  Their
only asymptotic difference is the constant $C_r$ absorbed into $R_r$.
\end{remark}

\subsection{The single periodic OEIS continuation}

A convenient one-function interpolation is
\begin{equation}
 \DO(x)=2^{x/2}\Gamma\!\left(\frac{x}{2}+1\right)
 \left(\frac2\pi\right)^{(1-\cos\pi x)/4}.
 \label{eq:OEIS-continuation}
\end{equation}
For even integers the periodic factor is $1$, while for odd integers it is
$\sqrt{2/\pi}$; hence $\DO(n)=n!!$ for all $n\ge0$.  This is the analytic
extension recorded in the OEIS entry \cite{OEIS}.

Let
\begin{equation}
 C_*=\frac14\log(2\pi),
 \qquad a=\frac14\log\frac\pi2>0.
 \label{eq:mean-periodic-constants}
\end{equation}
In the centred variable $s=x+1$, the exact relation developed below is
\begin{equation}
 \log\DO(s-1)
 =\frac{s}{2}(\log s-1)+C_*-a\cos(\pi s)+\mathcal C(s),
 \label{eq:OEIS-centered-exact}
\end{equation}
where $\mathcal C$ is the Borel-summed centred Stirling correction.  Thus the
periodic term remains $O(1)$ in the logarithm.  On inversion it produces an
$O(1/\log T)$ displacement, which is much larger than the
$O(1/(T\log T))$ first Bernoulli displacement.

The function $\DO$ is eventually strictly increasing.  Indeed,
\begin{equation}
 \frac{\DO'(x)}{\DO(x)}
 =\frac12\left(\log2+\psi\!\left(\frac{x}{2}+1\right)\right)
 +\frac{\pi}{4}\log\frac2\pi\,\sin(\pi x).
 \label{eq:OEIS-logder}
\end{equation}
The first term grows like $\tfrac12\log x$, while the second is bounded.
More explicitly, the standard bound
$\psi(z)>\log(z-1/2)$ for $z>1/2$ gives a lower bound
$\tfrac12\log(x+1)-\pi a$, which is positive for all sufficiently large $x$.
The eventual inverse $\DO_+^{-1}$ is therefore well defined; its complete
formal construction is given in \cref{sec:periodic-inverse}.

\subsection{The discrete generalized inverse}

For a fixed parity, define the largest same-parity index whose double
factorial does not exceed $y$ by
\begin{equation}
 J_r(y)=r+2\left\lfloor\frac{\D_r^{-1}(y)-r}{2}\right\rfloor,
 \qquad r\in\{0,1\},
 \label{eq:discrete-parity-inverse}
\end{equation}
whenever $y$ is in the relevant range.  The generalized inverse of the whole
sequence is then
\begin{equation}
 J(y)=\max\{J_0(y),J_1(y)\}.
 \label{eq:discrete-global-inverse}
\end{equation}
The smooth transseries computed later can be inserted into
\cref{eq:discrete-parity-inverse}.  Because the final floor is discontinuous,
a rigorous integer answer requires an error bound smaller than the distance to
the nearest parity threshold; this is discussed in \cref{sec:algorithm}.

\section{Direct asymptotics: two optimal normal forms}
\label{sec:direct}

Throughout, Bernoulli numbers are normalized by
\begin{equation}
 \frac{t}{\e^t-1}=\sum_{n\ge0}B_n\frac{t^n}{n!},
 \qquad B_1=-\frac12.
 \label{eq:bernoulli-convention}
\end{equation}
The classical Stirling series is
\begin{equation}
 \log\Gamma(z)
 \sim\left(z-\frac12\right)\log z-z+\frac12\log(2\pi)
 +\sum_{k\ge1}\frac{B_{2k}}{2k(2k-1)z^{2k-1}},
 \label{eq:stirling-log}
\end{equation}
valid in closed subsectors of $|\arg z|<\pi$ \cite{DLMFGamma}.

\subsection{Uncentred expansion}

\begin{theorem}[Uncentred logarithmic expansion]
\label{thm:uncentred-log}
As $x\to+\infty$,
\begin{equation}
 \log\D_r(x)
 \sim \frac{x}{2}(\log x-1)+\frac12\log x+C_r
 +\sum_{k\ge1}\frac{\delta_k}{x^{2k-1}},
 \label{eq:uncentred-log}
\end{equation}
where
\begin{equation}
 \delta_k=\frac{2^{2k-1}B_{2k}}{2k(2k-1)}
 =(-1)^{k-1}\frac{(2k-2)!\zeta(2k)}{\pi^{2k}}.
 \label{eq:delta-def}
\end{equation}
In particular,
\begin{equation}
 \delta_1=\frac16,\quad
 \delta_2=-\frac1{45},\quad
 \delta_3=\frac8{315},\quad
 \delta_4=-\frac8{105},\quad
 \delta_5=\frac{128}{297}.
 \label{eq:delta-first}
\end{equation}
\end{theorem}

\begin{proof}
Apply \cref{eq:stirling-log} to
$\Gamma(x/2+1)$ in the equivalent shifted form
\[
 \log\Gamma(z+1)
 \sim\left(z+\frac12\right)\log z-z+\frac12\log(2\pi)
 +\sum_{k\ge1}\frac{B_{2k}}{2k(2k-1)z^{2k-1}}
\]
with $z=x/2$, and add $(x/2)\log2+\log\kappa_r$.
The elementary terms reduce to the first three terms in
\cref{eq:uncentred-log}, while $z^{-(2k-1)}=2^{2k-1}x^{-(2k-1)}$.
Euler's evaluation of $\zeta(2k)$ gives the second form of $\delta_k$.
\end{proof}

Exponentiating gives
\begin{equation}
 \D_r(x)
 \sim A_r\sqrt{x}\left(\frac{x}{\e}\right)^{x/2}
 \sum_{n\ge0}\frac{g_n}{x^n},
 \qquad A_0=\sqrt\pi,\quad A_1=\sqrt2.
 \label{eq:uncentred-multiplicative}
\end{equation}
The coefficients $g_n$ are developed in \cref{sec:direct-coefficients}.

\subsection{The optimally centred expansion}

The shift $s=x+1$ removes the explicit square-root factor and annihilates all
even inverse powers in the logarithm.

\begin{theorem}[Centred logarithmic expansion]
\label{thm:centered-log}
Let $s=x+1$.  Then
\begin{equation}
 \log\D_r(s-1)
 \sim \frac{s}{2}(\log s-1)+C_r
 +\sum_{k\ge1}\frac{\beta_k}{s^{2k-1}},
 \label{eq:centered-log}
\end{equation}
where $\beta_k$ is given by \cref{eq:beta-def}.  The first terms are
\begin{equation}
 \beta_1=-\frac1{12},\quad
 \beta_2=\frac7{360},\quad
 \beta_3=-\frac{31}{1260},\quad
 \beta_4=\frac{127}{1680},\quad
 \beta_5=-\frac{511}{1188}.
 \label{eq:beta-first}
\end{equation}
Equivalently,
\begin{equation}
 \D_r(s-1)
 \sim A_r\left(\frac{s}{\e}\right)^{s/2}
 \sum_{n\ge0}\frac{c_n}{s^n}.
 \label{eq:centered-multiplicative}
\end{equation}
\end{theorem}

\begin{proof}
The duplication formula
\begin{equation}
 \Gamma(s)=\frac{2^{s-1}}{\sqrt\pi}
 \Gamma\!\left(\frac{s}{2}\right)
 \Gamma\!\left(\frac{s+1}{2}\right)
 \label{eq:duplication}
\end{equation}
(see, for example, \cite{DLMFDuplication})
rewrites \cref{eq:Dr-def} as
\begin{equation}
 \D_r(s-1)=\kappa_r\sqrt\pi\,2^{(1-s)/2}
 \frac{\Gamma(s)}{\Gamma(s/2)}.
 \label{eq:Dr-ratio}
\end{equation}
Subtracting the two Stirling expansions cancels the logarithmic square-root
term.  The $k$th correction is
\[
 \frac{B_{2k}}{2k(2k-1)}
 \left(s^{-(2k-1)}-(s/2)^{-(2k-1)}\right),
\]
which is precisely $\beta_ks^{-(2k-1)}$.  The constant is
$\log(\kappa_r\sqrt\pi)=C_r$.  Exponentiation gives
\cref{eq:centered-multiplicative}.
\end{proof}

At integer $n$, the two parity constants can be written in one expression:
\begin{equation}
 A_n=(2\pi)^{1/4}\left(\frac\pi2\right)^{(-1)^n/4}
 =\pi^{(1+(-1)^n)/4}2^{(1-(-1)^n)/4}.
 \label{eq:parity-amplitude}
\end{equation}
Hence the unified centred sequence expansion is
\begin{equation}
 n!!\sim A_n\left(\frac{n+1}{\e}\right)^{(n+1)/2}
 \sum_{j\ge0}\frac{c_j}{(n+1)^j}.
 \label{eq:sequence-centered-unified}
\end{equation}
The periodic amplitude is not an error term; it is part of the leading
transmonomial.

\section{Combinatorial formulae for the direct coefficients}
\label{sec:direct-coefficients}

The passage from the logarithmic series to the multiplicative series is an
instance of the exponential formula.  We record it in three equivalent forms:
a partition sum, a complete Bell polynomial, and a triangular recurrence.
This makes every direct coefficient computable without re-expanding an
exponential from scratch.

\subsection{Bell-polynomial conventions}

The exponential partial Bell polynomials are defined by
\begin{equation}
 \frac1{k!}\left(\sum_{m\ge1}x_m\frac{z^m}{m!}\right)^k
 =\sum_{n\ge k}\BellP{n}{k}(x_1,\ldots,x_{n-k+1})\frac{z^n}{n!},
 \label{eq:partial-bell-def}
\end{equation}
and the complete Bell polynomial is
\begin{equation}
 \BellC{n}(x_1,\ldots,x_n)
 =\sum_{k=0}^n\BellP{n}{k}(x_1,\ldots,x_{n-k+1}).
 \label{eq:complete-bell-def}
\end{equation}
Thus
\begin{equation}
 \exp\!\left(\sum_{m\ge1}x_m\frac{z^m}{m!}\right)
 =\sum_{n\ge0}\BellC{n}(x_1,\ldots,x_n)\frac{z^n}{n!}.
 \label{eq:bell-exponential}
\end{equation}
These conventions agree with the coefficient calculus used in
\cite{ProveItCCC,Comtet}.

\begin{theorem}[All direct multiplicative coefficients]
\label{thm:direct-coefficients}
Define sparse logarithmic coefficient sequences
\begin{equation}
 \ell^{(u)}_{2k-1}=\delta_k,\quad \ell^{(u)}_{2k}=0,
 \qquad
 \ell^{(c)}_{2k-1}=\beta_k,\quad \ell^{(c)}_{2k}=0.
 \label{eq:sparse-log-coeff}
\end{equation}
Then the coefficients in \cref{eq:uncentred-multiplicative} and
\cref{eq:centered-multiplicative} are
\begin{align}
 g_n&=\frac1{n!}\BellC{n}
 \bigl(1!\ell^{(u)}_1,2!\ell^{(u)}_2,\ldots,n!\ell^{(u)}_n\bigr),
 \label{eq:g-bell}\\
 c_n&=\frac1{n!}\BellC{n}
 \bigl(1!\ell^{(c)}_1,2!\ell^{(c)}_2,\ldots,n!\ell^{(c)}_n\bigr).
 \label{eq:c-bell}
\end{align}
Equivalently,
\begin{align}
 g_n&=\sum_{\substack{m_1,m_2,\ldots\ge0\\
             \sum_{k\ge1}(2k-1)m_k=n}}
       \prod_{k\ge1}\frac{\delta_k^{m_k}}{m_k!},
 \label{eq:g-partition}\\
 c_n&=\sum_{\substack{m_1,m_2,\ldots\ge0\\
             \sum_{k\ge1}(2k-1)m_k=n}}
       \prod_{k\ge1}\frac{\beta_k^{m_k}}{m_k!}.
 \label{eq:c-partition}
\end{align}
They obey the recurrences
\begin{align}
 g_0&=1,&
 n g_n&=\sum_{1\le k\le(n+1)/2}(2k-1)\delta_k\,g_{n-2k+1},
 \label{eq:g-recurrence}\\
 c_0&=1,&
 n c_n&=\sum_{1\le k\le(n+1)/2}(2k-1)\beta_k\,c_{n-2k+1}.
 \label{eq:c-recurrence}
\end{align}
\end{theorem}

\begin{proof}
Put $z=x^{-1}$ in the uncentred expansion.  The correction factor is
\[
 \exp\!\left(\sum_{k\ge1}\delta_kz^{2k-1}\right).
\]
The formula in \cref{eq:g-bell} follows from \cref{eq:bell-exponential}, and expanding
the exponential as a product of
$\sum_{m_k\ge0}\delta_k^{m_k}z^{(2k-1)m_k}/m_k!$ gives
\cref{eq:g-partition}.  If $G(z)=\sum_{n\ge0}g_nz^n$, logarithmic
differentiation gives
\[
 zG'(z)=G(z)\sum_{k\ge1}(2k-1)\delta_kz^{2k-1}.
\]
Coefficient extraction proves \cref{eq:g-recurrence}.  Replacing $\delta_k$
by $\beta_k$ proves all centred statements.
\end{proof}

The first coefficients are as follows.
\begin{table}[htbp]
\centering
\caption{Multiplicative coefficients in the uncentred and centred normal forms.}
\label{tab:direct-coefficients}
\small
\begin{tabular}{c@{\qquad}r@{\qquad}r}
\toprule
$n$ & $g_n$ & $c_n$\\
\midrule
0 & $1$ & $1$\\
1 & $1/6$ & $-1/12$\\
2 & $1/72$ & $1/288$\\
3 & $-139/6480$ & $1003/51840$\\
4 & $-571/155520$ & $-4027/2488320$\\
5 & $163879/6531840$ & $-5128423/209018880$\\
6 & $5246819/1175731200$ & $168359651/75246796800$\\
7 & $-534703531/7054387200$ & $68168266699/902961561600$\\
8 & $-4483131259/338610585600$ & $-587283555451/86684309913600$\\
\bottomrule
\end{tabular}
\end{table}

\begin{remark}[Arithmetic]
All $g_n$ and $c_n$ are rational because the logarithmic coefficients are
rational.  The Bell formula is preferable for structural arguments, while the
recurrence is normally preferable for exact computation: it produces the first
$N$ coefficients with $O(N^2)$ rational operations and exposes the sparse odd
support directly.
\end{remark}

\section{Exact Borel sums, remainder bounds, and large order}
\label{sec:borel}

The Stirling series is divergent, but in the present problem its Borel
transform is elementary.  This gives more information than a formal expansion:
it identifies the exact analytic function represented by the series and makes
its remainder mechanism transparent.

\subsection{Borel convention and Binet's remainder}

For a formal inverse-power series
\begin{equation}
 F(s)\sim\sum_{m\ge0}\frac{a_m}{s^{m+1}},
 \label{eq:borel-convention-series}
\end{equation}
we use the formal Borel transform
\begin{equation}
 \widehat F(t)=\sum_{m\ge0}a_m\frac{t^m}{m!}.
 \label{eq:borel-convention-transform}
\end{equation}
Whenever the integral exists along the positive ray, its Laplace sum is
\begin{equation}
 \mathscr S_0F(s)=\int_0^\infty\e^{-st}\widehat F(t)\dd t.
 \label{eq:borel-laplace}
\end{equation}

For $\Re z>0$, define the exact Binet correction
\begin{equation}
 \Omega(z)=\log\Gamma(z)-\left(z-\frac12\right)\log z+z
 -\frac12\log(2\pi).
 \label{eq:Omega-def}
\end{equation}
Binet's integral can be written as
\begin{equation}
 \Omega(z)=\int_0^\infty\e^{-zt}
 \left(\frac{\coth(t/2)}{2t}-\frac1{t^2}\right)\dd t.
 \label{eq:Binet-Laplace}
\end{equation}
The apparent singularity at $t=0$ is removable.  Taylor expansion of the
kernel gives Stirling's Bernoulli series, while exponential decay at infinity
makes the integral convergent for $\Re z>0$.

\subsection{The two elementary Borel transforms}

Define the exact uncentred and centred logarithmic corrections by
\begin{align}
 \mathcal U(x)&=\log\D_r(x)
 -\frac{x}{2}(\log x-1)-\frac12\log x-C_r,
 \label{eq:U-def}\\
 \mathcal C(s)&=\log\D_r(s-1)
 -\frac{s}{2}(\log s-1)-C_r.
 \label{eq:C-def}
\end{align}
They do not depend on $r$.

\begin{theorem}[Exact Borel--Laplace representations]
\label{thm:exact-borel}
For $x>0$ and $s>0$,
\begin{align}
 \mathcal U(x)
 &=\int_0^\infty\e^{-xt}
 \left(\frac{\coth t}{2t}-\frac1{2t^2}\right)\dd t,
 \label{eq:U-Laplace}\\
 \mathcal C(s)
 &=\int_0^\infty\e^{-st}
 \left(\frac{1}{2t\sinh t}-\frac1{2t^2}\right)\dd t.
 \label{eq:C-Laplace}
\end{align}
Their Borel transforms admit the Mittag--Leffler sums
\begin{align}
 \widehat{\mathcal U}(t)
 &=\frac{\coth t}{2t}-\frac1{2t^2}
 =\sum_{m=1}^\infty\frac1{\pi^2m^2+t^2},
 \label{eq:U-partial-fractions}\\
 \widehat{\mathcal C}(t)
 &=\frac{1}{2t\sinh t}-\frac1{2t^2}
 =-\sum_{m=1}^\infty\frac{(-1)^{m-1}}{\pi^2m^2+t^2}.
 \label{eq:C-partial-fractions}
\end{align}
Consequently \cref{eq:uncentred-log,eq:centered-log} are Borel summable along
the positive real direction, and their Borel sums are the exact gamma
expressions.
\end{theorem}

\begin{proof}
The uncentred correction is $\mathcal U(x)=\Omega(x/2)$.  Substitute $t=2u$
in \cref{eq:Binet-Laplace} to obtain \cref{eq:U-Laplace}.  From
\cref{eq:Dr-ratio}, the centred correction is
$\mathcal C(s)=\Omega(s)-\Omega(s/2)$.  Subtracting the corresponding kernels
and using
\[
 \coth(t/2)-\coth t=\frac1{\sinh t}
\]
gives \cref{eq:C-Laplace}.  The classical Mittag--Leffler expansion
\[
 \frac{\coth t}{t}=\frac1{t^2}
 +2\sum_{m\ge1}\frac1{t^2+\pi^2m^2}
\]
proves \cref{eq:U-partial-fractions}.  Splitting this sum into even and odd
indices, or equivalently subtracting the scaled kernel, gives
\cref{eq:C-partial-fractions}.  Both kernels are analytic on the positive
integration ray and have at most exponential growth in sectors avoiding their
poles, so the Laplace integrals are unambiguous there.
\end{proof}

Expanding the partial fractions at $t=0$ immediately recovers
\begin{align}
 \widehat{\mathcal U}(t)
 &=\sum_{k\ge1}(-1)^{k-1}
   \frac{\zeta(2k)}{\pi^{2k}}t^{2k-2},
 \label{eq:U-borel-series}\\
 \widehat{\mathcal C}(t)
 &=\sum_{k\ge1}(-1)^k
   \frac{\etaD(2k)}{\pi^{2k}}t^{2k-2},
 \label{eq:C-borel-series}
\end{align}
which is equivalent to \cref{eq:delta-def,eq:beta-def} after Laplace
integration.

\subsection{Sharp alternating remainder bounds}

\begin{theorem}[First-omitted-term bounds]
\label{thm:remainder-bounds}
For $x>0$, $s>0$, and $N\ge0$, put
\begin{align}
 U_N(x)&=\sum_{k=1}^{N}\frac{\delta_k}{x^{2k-1}},
 \label{eq:UN-def}\\
 C_N(s)&=\sum_{k=1}^{N}\frac{\beta_k}{s^{2k-1}}.
 \label{eq:CN-def}
\end{align}
Then
\begin{align}
 0&<(-1)^N\bigl(\mathcal U(x)-U_N(x)\bigr)
 <\frac{|\delta_{N+1}|}{x^{2N+1}},
 \label{eq:U-remainder-bound}\\
 0&<(-1)^{N+1}\bigl(\mathcal C(s)-C_N(s)\bigr)
 <\frac{|\beta_{N+1}|}{s^{2N+1}}.
 \label{eq:C-remainder-bound}
\end{align}
The convention is that the empty sum for $N=0$ equals zero.
\end{theorem}

\begin{proof}
For $a>0$,
\begin{equation}
 \frac1{a^2+t^2}
 =\sum_{j=0}^{N-1}\frac{(-1)^jt^{2j}}{a^{2j+2}}
 +\frac{(-1)^Nt^{2N}}{a^{2N}(a^2+t^2)}.
 \label{eq:partial-fraction-remainder}
\end{equation}
Insert this identity into \cref{eq:U-partial-fractions}, sum over
$a=\pi m$, and integrate.  The remainder has sign $(-1)^N$; replacing
$a^2+t^2$ by $a^2$ gives the upper bound, and
$\int_0^\infty\e^{-xt}t^{2N}\dd t=(2N)!x^{-(2N+1)}$ yields
\cref{eq:U-remainder-bound}.

For the centred transform, the residual kernel is
\[
 (-1)^{N+1}t^{2N}S_N(t),\qquad
 S_N(t)=\sum_{m\ge1}\frac{(-1)^{m-1}}
 { (\pi m)^{2N}\bigl((\pi m)^2+t^2\bigr)}.
\]
To see that $S_N(t)>0$ and decreases with $t^2$, use
\[
 S_N(t)=\int_0^\infty\e^{-t^2v}
 \sum_{m\ge1}\frac{(-1)^{m-1}\e^{-\pi^2m^2v}}{(\pi m)^{2N}}\dd v.
\]
For every $v\ge0$ the inner series is an alternating series with strictly
decreasing positive terms, hence is positive.  Therefore
$0<S_N(t)\le S_N(0)=\etaD(2N+2)/\pi^{2N+2}$.  Laplace integration now gives
\cref{eq:C-remainder-bound}.
\end{proof}

These inequalities are particularly useful for inversion: they give a
certified logarithmic residual without evaluating a huge double factorial.
They also show directly that the centred series alternates with its first term
negative.

\subsection{Large-order growth and optimal truncation}

The exact coefficient formulae imply
\begin{align}
 \delta_k&=(-1)^{k-1}\frac{(2k-2)!}{\pi^{2k}}
 \left(1+O(2^{-2k})\right),
 \label{eq:delta-large-order}\\
 \beta_k&=(-1)^k\frac{(2k-2)!}{\pi^{2k}}
 \left(1+O(2^{-2k})\right).
 \label{eq:beta-large-order}
\end{align}
Thus the $k$th centred term is smallest when
$2k\approx\pi s$.  Stirling's formula then gives the least-term envelope
\begin{equation}
 \min_k\left|\frac{\beta_k}{s^{2k-1}}\right|
 =O\!\left(s^{-1/2}\e^{-\pi s}\right).
 \label{eq:least-term-envelope}
\end{equation}
The nearest Borel singularities are the simple poles $t=\pm\ii\pi$; the full
lattice is $t=\pm\ii\pi m$, $m\ge1$.  This pole geometry simultaneously
explains the factorial coefficient growth and the scale
\cref{eq:least-term-envelope}.

\begin{remark}[Exponentially improved interpretation]
The optimal remainder on the positive real axis has the envelope
$\e^{-\pi s}$, but the analytically continued hyperasymptotic sectors are not
plain real exponentials.  They are terminant-weighted contributions associated
with singulants $\pm\ii\pi ms$.  Away from a Stokes direction the terminant
multiplier itself is exponentially small.  This distinction matters: appending
$\e^{-\pi s}$ with an arbitrary constant to the divergent Bernoulli series is
not a valid transseries completion.  The exact Borel integral
\cref{eq:C-Laplace} is the unambiguous positive-real completion.
\end{remark}

\section{Lambert normalization of the branchwise inverse}
\label{sec:lambert-core}

The centred phase is ideal for inversion because its dominant part has an
exact Lambert-$W$ inverse and its perturbation occupies only odd inverse powers.
We first isolate the core, then solve the perturbation on a well-ordered
$q$-grid.

\subsection{Exact inversion of the core}

Let
\begin{equation}
 \Phi_0(s)=s(\log s-1),\qquad s>1.
 \label{eq:core-phase}
\end{equation}
For $R>0$, the equation $\Phi_0(T)=2R$ has an exact solution.

\begin{lemma}[Lambert core]
\label{lem:lambert-core}
Define
\begin{equation}
 w=\W_0(2R/\e),
 \qquad T=\frac{2R}{w}=\e^{w+1},
 \qquad \Lambda=w+1=\log T.
 \label{eq:lambert-core}
\end{equation}
Then
\begin{equation}
 T(\log T-1)=2R,
 \qquad \Phi_0'(T)=\Lambda,
 \qquad \frac{\dd T}{\dd R}=\frac{2}{\Lambda}.
 \label{eq:core-identities}
\end{equation}
\end{lemma}

\begin{proof}
Since $w\e^w=2R/\e$, one has
$T=2R/w=\e^{w+1}$ and therefore $\log T-1=w$.  Multiplication gives the first
identity.  The other two follow by differentiation.
\end{proof}

For the parity branch $r$, take $R=R_r=\log y-C_r$.  The leading inverse is
$x=T_r-1$.  Retaining $W_0$ exactly packages infinitely many outer logarithmic
terms into one analytic quantity; this is both more accurate and algebraically
cleaner than expanding $W_0$ at the outset.

\subsection{The exact normalized formal equation}

Multiply \cref{eq:centered-log} by $2$ and use $\alpha_k=2\beta_k$.  Formally,
solving $\D_r(s-1)=y$ is equivalent to
\begin{equation}
 2R_r=s(\log s-1)+\sum_{k\ge1}\frac{\alpha_k}{s^{2k-1}}.
 \label{eq:inverse-forward-formal}
\end{equation}
Let $T=T_r$, $\Lambda=\Lambda_r$, $q=T^{-1}$, and write
\begin{equation}
 s=T(1+E),
 \qquad E=O(q^2).
 \label{eq:relative-displacement}
\end{equation}
Subtracting the core equation and dividing by $T$ gives
\begin{equation}
 \Lambda E+\varphi(E)
 +\sum_{k\ge1}\alpha_kq^{2k}(1+E)^{-(2k-1)}=0,
 \label{eq:normalized-inverse-equation}
\end{equation}
where
\begin{equation}
 \varphi(u)=(1+u)\log(1+u)-u
 =\sum_{m\ge2}\frac{(-1)^m}{m(m-1)}u^m.
 \label{eq:varphi-def}
\end{equation}
The support is $2\N$: if $E$ contains only even powers of $q$, so does the
right-hand side, and the linear coefficient $\Lambda$ is invertible in the
coefficient field $\Q(\Lambda)$.

\section{All-orders parity-resolved inverse transseries}
\label{sec:branch-inverse}

\subsection{Triangular recurrence and analytic remainder}

\begin{theorem}[Complete branchwise inverse expansion]
\label{thm:branch-inverse}
Let $r\in\{0,1\}$ and define $R_r,w_r,T_r,\Lambda_r$ by
\cref{eq:inverse-core-summary}.  For every fixed $N\ge0$,
\begin{equation}
 \D_r^{-1}(y)
 =T_r-1+\sum_{n=1}^{N}
 \frac{d_{2n-1}(\Lambda_r)}{T_r^{2n-1}}
 +O\!\left(\frac{1}{T_r^{2N+1}\Lambda_r}\right),
 \qquad y\to+\infty.
 \label{eq:branch-inverse-all-orders}
\end{equation}
The coefficients are uniquely determined as follows.  Put
\begin{equation}
 E_{n-1}(q)=\sum_{j=1}^{n-1}e_j(\Lambda)q^{2j}.
 \label{eq:E-truncated}
\end{equation}
Then
\begin{equation}
 e_n(\Lambda)=-\frac1\Lambda\coeff{q}{2n}
 \left[
  \varphi(E_{n-1})+
  \sum_{k=1}^{n}\alpha_kq^{2k}
       (1+E_{n-1})^{-(2k-1)}
 \right],
 \label{eq:inverse-triangular-recurrence}
\end{equation}
with
\begin{equation}
 d_{2n-1}(\Lambda)=e_n(\Lambda).
 \label{eq:d-e-identification}
\end{equation}
Furthermore,
\begin{equation}
 d_{2n-1}(\Lambda)\in\Lambda^{-1}\Q[\Lambda^{-1}],
 \qquad
 \deg_{\Lambda^{-1}}d_{2n-1}\le2n-1,
 \label{eq:d-ring-degree}
\end{equation}
and
\begin{equation}
 d_{2n-1}(\Lambda)=-\frac{\alpha_n}{\Lambda}
 +O(\Lambda^{-2}).
 \label{eq:d-leading-slope}
\end{equation}
\end{theorem}

\begin{proof}
Seek a solution of \cref{eq:normalized-inverse-equation} in the form
$E=\sum_{n\ge1}e_nq^{2n}$.  At order $q^{2n}$, the only term containing the
unknown $e_n$ is $\Lambda e_nq^{2n}$.  The nonlinear series $\varphi(E)$
starts with $E^2/2$, and every term in the Bernoulli sector has positive
$q$-valuation.  Consequently the remaining coefficient of $q^{2n}$ depends
only on $e_1,\ldots,e_{n-1}$, and solving the coefficient equation gives
\cref{eq:inverse-triangular-recurrence}.  This proves formal existence and
uniqueness by induction.

The same induction proves \cref{eq:d-ring-degree}.  The coefficient inside the
brackets is a polynomial in $\Lambda^{-1}$ of degree at most $2n-2$; division
by $\Lambda$ gives the claimed bound.  The only contribution at $q^{2n}$ that
is independent of lower $e_j$ is $\alpha_nq^{2n}$, proving
\cref{eq:d-leading-slope}.

For the analytic statement, truncate the exact centred Borel sum after the
term $\beta_{N+1}s^{-(2N+1)}$.  By \cref{thm:remainder-bounds}, the omitted
forward logarithmic residual is $O(s^{-(2N+3)})$.  Substituting the formal
inverse through $e_Nq^{2N}$ cancels the doubled phase through order $q^{2N}$
and leaves a residual $O(T^{-(2N+1)})$.  The derivative of the doubled phase
is
\begin{equation}
 \log s-2\sum_{k\ge1}(2k-1)\beta_ks^{-2k}
 =\Lambda+O(T^{-2}).
 \label{eq:doubled-phase-slope}
\end{equation}
For large $T$ it is bounded below by a fixed positive multiple of $\Lambda$.
The mean-value theorem therefore transfers the forward residual to an inverse
error $O(T^{-(2N+1)}/\Lambda)$.
\end{proof}

\begin{remark}[Uniformity in parity]
The implied constants in \cref{eq:branch-inverse-all-orders} can be chosen
uniformly for $r=0,1$.  The coefficient recurrence contains no parity label;
$r$ enters only through the bounded constants $C_r$ in $R_r$.
\end{remark}

\subsection{Explicit inverse blocks}

Successive use of \cref{eq:inverse-triangular-recurrence} gives
\begin{align}
 d_1(\Lambda)={}&\frac1{6\Lambda},
 \label{eq:d1}\\[0.35em]
 d_3(\Lambda)={}&-\frac{14\Lambda^2+10\Lambda+5}
 {360\Lambda^3},
 \label{eq:d3}\\[0.35em]
 d_5(\Lambda)={}&\frac{1}{45360\Lambda^5}
 \bigl(2232\Lambda^4+1176\Lambda^3+714\Lambda^2
 +350\Lambda+105\bigr),
 \label{eq:d5}\\[0.35em]
 d_7(\Lambda)={}&-\frac{1}{5443200\Lambda^7}
 \bigl(822960\Lambda^6+292536\Lambda^5+136956\Lambda^4
 \notag\\[-0.2em]
 &\hspace{5.1em}{}+68040\Lambda^3+32970\Lambda^2
 +12250\Lambda+2625\bigr),
 \label{eq:d7}\\[0.35em]
 d_9(\Lambda)={}&\frac{1}{359251200\Lambda^9}
 \bigl(309052800\Lambda^8+77920128\Lambda^7+26024592\Lambda^6
 \notag\\[-0.2em]
 &\hspace{2.8em}{}+10019592\Lambda^5+4516028\Lambda^4
 +2023560\Lambda^3+812350\Lambda^2
 \notag\\[-0.2em]
 &\hspace{2.8em}{}+242550\Lambda+40425\bigr).
 \label{eq:d9}
\end{align}
Thus the first three correction blocks are
\begin{align}
 \D_r^{-1}(y)={}&T_r-1+\frac{1}{6T_r\Lambda_r}
 -\frac{14\Lambda_r^2+10\Lambda_r+5}
 {360T_r^3\Lambda_r^3}
 \notag\\
 &+\frac{2232\Lambda_r^4+1176\Lambda_r^3+714\Lambda_r^2
 +350\Lambda_r+105}{45360T_r^5\Lambda_r^5}
 +O\!\left(\frac1{T_r^7\Lambda_r}\right).
 \label{eq:first-three-inverse-blocks}
\end{align}

For completeness, the next block is
\begin{align}
 d_{11}(\Lambda)={}&-\frac{1}{5884534656000\Lambda^{11}}
 \bigl(46195687238400\Lambda^{10}+8854370366400\Lambda^9
 \notag\\[-0.2em]
 &+2171643318000\Lambda^8+605515022112\Lambda^7
 +212869408752\Lambda^6+84136852800\Lambda^5
 \notag\\[-0.2em]
 &+35589153600\Lambda^4+14234019800\Lambda^3
 +4868113250\Lambda^2+1213962750\Lambda
 \notag\\[-0.2em]
 &+165540375\bigr).
 \label{eq:d11}
\end{align}
The rapidly growing numerators are an argument for using the general formulae
rather than storing a table.

\subsection{Low-order derivation by hand}

At order $q^2$, \cref{eq:normalized-inverse-equation} gives
$\Lambda e_1+\alpha_1=0$.  Since $\alpha_1=-1/6$,
\[
 e_1=\frac1{6\Lambda}.
\]
At order $q^4$ one obtains
\begin{equation}
 \Lambda e_2+\frac12e_1^2-\alpha_1e_1+\alpha_2=0.
 \label{eq:e2-hand}
\end{equation}
Substitution of $\alpha_1=-1/6$, $\alpha_2=7/180$ yields
\cref{eq:d3}.  At every subsequent order the calculation is finite; the
number of additive partitions grows, but no new analytic idea is required.

\section{A finite closed formula for every inverse coefficient}
\label{sec:closed-coefficients}

The triangular recurrence is optimal for computation, but it conceals the
combinatorial content.  Lagrange--B\"urmann inversion converts it into a finite,
non-recursive expression involving two coefficient arrays: one built from the
Bernoulli data and one from the universal nonlinear kernel $\varphi$.

\subsection{Bernoulli convolution and kernel arrays}

Define
\begin{equation}
 A(z)=\sum_{k\ge1}\alpha_kz^k,
 \qquad
 A_{n,p}=\coeff{z}{n}A(z)^p.
 \label{eq:A-np-def}
\end{equation}
By the defining identity for partial Bell polynomials,
\begin{equation}
 A_{n,p}=\frac{p!}{n!}
 \BellP{n}{p}\bigl(1!\alpha_1,2!\alpha_2,\
 \ldots,(n-p+1)!\alpha_{n-p+1}\bigr).
 \label{eq:A-np-bell}
\end{equation}
For integers $n,j,p$ define
\begin{equation}
 K_{n;j,p}=\coeff{u}{j-1}
 \varphi(u)^{j-p}(1+u)^{-(2n-p)}.
 \label{eq:K-def}
\end{equation}
If $p>j$, or if $2(j-p)>j-1$, this coefficient is zero.

To express $K_{n;j,p}$ itself through Bell polynomials, put
\begin{equation}
 \nu_1=0,
 \qquad \nu_m=(-1)^m(m-2)!\quad(m\ge2),
 \label{eq:nu-def}
\end{equation}
so that $\varphi(u)=\sum_{m\ge1}\nu_mu^m/m!$.  With
$r=j-p$ and the conventions $\BellP{0}{0}=1$ and
$\BellP{h}{0}=0$ for $h>0$, one has the finite formula
\begin{align}
 K_{n;j,p}={}&
 \sum_{h=2r}^{j-1}
 \frac{r!}{h!}\BellP{h}{r}(\nu_1,\ldots,\nu_{h-r+1})
 (-1)^{j-1-h}
 \notag\\[-0.2em]
 &\hspace{5em}{}\times
 \binom{2n-p+j-2-h}{j-1-h}.
 \label{eq:K-bell}
\end{align}
When $r=0$, only the $h=0$ term is retained.

\begin{theorem}[Closed coefficient formula]
\label{thm:closed-coefficients}
Write
\begin{equation}
 d_{2n-1}(\Lambda)=\sum_{j=1}^{2n-1}c_{n,j}\Lambda^{-j}.
 \label{eq:c-nj-expansion}
\end{equation}
Then, for $1\le j\le2n-1$,
\begin{equation}
 \boxed{
 c_{n,j}=\frac{(-1)^j}{j}
 \sum_{p=\lceil(j+1)/2\rceil}^{\min(j,n)}
 \binom{j}{p}A_{n,p}K_{n;j,p}.}
 \label{eq:closed-c-nj}
\end{equation}
\Cref{eq:A-np-bell,eq:K-bell} make the right-hand side a completely
finite expression in Bernoulli numbers, binomial coefficients, and partial
Bell polynomials.
\end{theorem}

\begin{proof}
Introduce an auxiliary homotopy parameter $\zeta$ and rewrite
\cref{eq:normalized-inverse-equation} as
\begin{equation}
 E=-\frac{\zeta}{\Lambda}H(q,E),
 \qquad
 H(q,u)=\varphi(u)+
 \sum_{k\ge1}\alpha_kq^{2k}(1+u)^{-(2k-1)}.
 \label{eq:homotopy-equation}
\end{equation}
The formal Lagrange--B\"urmann formula gives
\begin{equation}
 E(q,\zeta)=\sum_{j\ge1}\frac{(-1)^j\zeta^j}{j!\Lambda^j}
 \left.\frac{\dd^{j-1}}{\dd u^{j-1}}H(q,u)^j\right|_{u=0}.
 \label{eq:LB-homotopy}
\end{equation}
Set $\zeta=1$ and extract $q^{2n}$.  In $H^j$, choose $p$ of the $j$
factors from the Bernoulli sector and the other $j-p$ from $\varphi$.
The binomial choice contributes $\binom jp$.  If the selected Bernoulli
indices sum to $n$, then
\[
 \prod_{\ell=1}^p(1+u)^{-(2k_\ell-1)}
 =(1+u)^{-(2n-p)},
\]
and their coefficient sum is $A_{n,p}$.  The remaining coefficient of
$u^{j-1}$ is $K_{n;j,p}$.  Since differentiation contributes $(j-1)!$, the
factor $1/j!$ becomes $1/j$, proving \cref{eq:closed-c-nj}.

The factor $\varphi^{j-p}$ starts in degree $2(j-p)$, so a nonzero
$u^{j-1}$ coefficient requires $p\ge\lceil(j+1)/2\rceil$.  As $p\le n$, this
also forces $j\le2n-1$, proving directly that the sum in
\cref{eq:c-nj-expansion} terminates.  Finally, expanding $\varphi^r$ by
\cref{eq:partial-bell-def} and multiplying by the binomial series for
$(1+u)^{-(2n-p)}$ gives \cref{eq:K-bell}.
\end{proof}

\begin{methodbox}
There are now two all-orders algorithms.
\begin{itemize}
\item The recurrence in \cref{eq:inverse-triangular-recurrence}
      generates complete blocks $d_1,d_3,\ldots,d_{2N-1}$ with aggressive
      truncation in $q$ and is the fastest route in a computer algebra system.
\item The closed formula in \cref{eq:closed-c-nj} computes any isolated scalar
      coefficient $c_{n,j}$ without constructing lower inverse blocks.  It is
      the better formula for sign, support, denominator, and boundary analyses.
\end{itemize}
Both use exactly the Bell-polynomial/Fa\`a di Bruno coefficient machinery of
combinatorial coefficient calculus.
\end{methodbox}

\subsection{Boundary coefficients}

Two edges of the triangular array $c_{n,j}$ simplify immediately.

\begin{corollary}[First and deepest inverse-logarithmic coefficients]
\label{cor:boundary-coefficients}
For $n\ge1$,
\begin{align}
 c_{n,1}&=-\alpha_n,
 \label{eq:c-n1}\\
 c_{n,2n-1}&=-\frac1{2^{n-1}(2n-1)}
 \binom{2n-1}{n}\alpha_1^n.
 \label{eq:c-deepest}
\end{align}
\end{corollary}

\begin{proof}
For $j=1$, only $p=1$ contributes, $K_{n;1,1}=1$, and
$A_{n,1}=\alpha_n$.  For $j=2n-1$, the support restriction forces $p=n$.
Then $A_{n,n}=\alpha_1^n$, and the required coefficient is the lowest term of
$\varphi(u)^{n-1}$, namely
$[u^{2n-2}](u^2/2)^{n-1}=2^{1-n}$.  Substitution in
\cref{eq:closed-c-nj} proves the result.
\end{proof}

\subsection{Global alternating sign law}

The signs visible in \cref{eq:d1,eq:d3,eq:d5,eq:d7,eq:d9} persist in every
coefficient of every block.

\begin{theorem}[Coefficientwise sign theorem]
\label{thm:sign-law}
For every $n\ge1$,
\begin{equation}
 (-1)^{n+1}d_{2n-1}(\Lambda)
 \in\Lambda^{-1}\Q_{>0}[\Lambda^{-1}].
 \label{eq:sign-law}
\end{equation}
In particular, every nonzero scalar coefficient $c_{n,j}$ has sign
$(-1)^{n+1}$.
\end{theorem}

\begin{proof}
By \cref{eq:alpha-def}, write
\begin{equation}
 \alpha_k=(-1)^k a_k^+,
 \qquad
 a_k^+=\frac{2(2k-2)!\etaD(2k)}{\pi^{2k}}>0.
 \label{eq:alpha-positive}
\end{equation}
Set $z=-q^2$ and $E=-P$.  Since
\begin{equation}
 \varphi(-P)=\sum_{m\ge2}\frac{P^m}{m(m-1)},
 \label{eq:varphi-positive}
\end{equation}
all coefficients are positive, and the normalized equation becomes
\begin{equation}
 \Lambda P=
 \sum_{m\ge2}\frac{P^m}{m(m-1)}
 +\sum_{k\ge1}a_k^+z^k(1-P)^{-(2k-1)}.
 \label{eq:positive-fixed-point}
\end{equation}
Starting from $P=0$ and solving coefficient by coefficient first gives
$P=\sum_{n\ge1}p_n(\Lambda^{-1})z^n$ with nonnegative rational
coefficients.  Strict positivity at every power
$\Lambda^{-j}$, $1\le j\le2n-1$, follows constructively.  For
$1\le j\le n$, take the forcing term with index $n-j+1$ and the
$P^{j-1}$ term in its positive binomial expansion, using the leading
$a_1^+z/\Lambda$ contribution in every copy of $P$.  For
$n<j\le2n-1$, use a positive rooted-tree contribution with $n$ leaves and
$j-n$ internal vertices from the nonlinear series
$\sum_{m\ge2}P^m/(m(m-1))$; such a tree exists for every
$1\le j-n\le n-1$.  Hence
$p_n\in\Lambda^{-1}\Q_{>0}[\Lambda^{-1}]$.  Returning to $q$,
\[
 E(q)=-P(-q^2)=\sum_{n\ge1}(-1)^{n+1}p_nq^{2n}.
\]
Thus $d_{2n-1}=(-1)^{n+1}p_n$, proving the theorem.
\end{proof}

\subsection{Scaling relation with the inverse-Gamma blocks}

Let $a_k^{\Gamma}=B_{2k}(1/2)/(2k(2k-1))$ be the half-shifted logarithmic
coefficients used for inverse Gamma.  Since
\begin{equation}
 \alpha_k=4^k a_k^{\Gamma},
 \label{eq:alpha-gamma-scale}
\end{equation}
the normalized double-factorial equation is obtained from the normalized
inverse-Gamma equation by replacing $q$ with $2q$, while leaving $\Lambda$
formal.  Therefore
\begin{equation}
 d_{2n-1}^{!!}(\Lambda)=4^n d_{2n-1}^{\Gamma}(\Lambda).
 \label{eq:inverse-gamma-scaling}
\end{equation}
This identity is an independent all-orders check on both recurrences.  It is a
coefficient-polynomial identity; the numerical Lambert variables used in the
two inverse problems need not coincide.

\section{Outer iterated-logarithmic expansion}
\label{sec:outer}

The Lambert form is the natural compact transseries.  When an expansion using
only elementary iterated logarithms is required, $W$ can itself be expanded.
The resulting coefficient calculus is again governed by classical
combinatorial numbers.

\subsection{Unsigned-Stirling formula for Lambert \texorpdfstring{$W$}{W}}

Set
\begin{equation}
 Z=\frac{2R}{\e},
 \qquad L_1=\log Z,
 \qquad L_2=\log L_1.
 \label{eq:L1-L2-def}
\end{equation}
For $Z\to+\infty$, the complete Poincar\'e expansion is
\begin{equation}
 \W_0(Z)\sim L_1-L_2+
 \sum_{a\ge0}\sum_{b\ge1}
 \frac{(-1)^a}{b!}
 \stirlingone{a+b}{a+1}
 \frac{L_2^b}{L_1^{a+b}},
 \label{eq:W-Stirling-expansion}
\end{equation}
where $\stirlingone{n}{k}$ is the unsigned Stirling number of the first
kind \cite{Corless,DLMFW}.  This double sum is understood by increasing total
power $a+b$.

Write
\begin{equation}
 \frac{\W_0(Z)}{L_1}=1+\sum_{m\ge1}\frac{u_m(L_2)}{L_1^m}.
 \label{eq:u-def}
\end{equation}
Then
\begin{align}
 u_1(L_2)&=-L_2,
 \label{eq:u1}\\
 u_m(L_2)&=
 \sum_{b=1}^{m-1}\frac{(-1)^{m-1-b}}{b!}
 \stirlingone{m-1}{m-b}L_2^b,
 \qquad m\ge2.
 \label{eq:um-general}
\end{align}
Thus every outer coefficient is an explicit polynomial in $L_2$ with integer
Stirling data.

\subsection{Reciprocal coefficients and the core}

Define $\tau_m(L_2)$ by
\begin{equation}
 \frac1{\W_0(Z)}\sim\frac1{L_1}
 \sum_{m\ge0}\frac{\tau_m(L_2)}{L_1^m}.
 \label{eq:tau-def}
\end{equation}
Then
\begin{equation}
 \tau_0=1,
 \qquad
 \tau_m=-\sum_{j=1}^m u_j\tau_{m-j}.
 \label{eq:tau-recurrence}
\end{equation}
A non-recursive partition form is
\begin{equation}
 \tau_m=
 \sum_{\substack{k_1,\ldots,k_m\ge0\\
                  \sum_{j=1}^mjk_j=m}}
 (-1)^K\frac{K!}{k_1!\cdots k_m!}
 \prod_{j=1}^m u_j^{k_j},
 \qquad K=\sum_{j=1}^m k_j.
 \label{eq:tau-partition}
\end{equation}
The first polynomials are
\begin{align}
 \tau_0={}&1,
 \notag\\
 \tau_1={}&L_2,
 \notag\\
 \tau_2={}&L_2(L_2-1),
 \notag\\
 \tau_3={}&\frac12L_2(2L_2^2-5L_2+2),
 \notag\\
 \tau_4={}&\frac16L_2(6L_2^3-26L_2^2+27L_2-6),
 \label{eq:tau-first}\\
 \tau_5={}&\frac1{12}L_2
 (12L_2^4-77L_2^3+142L_2^2-84L_2+12).
 \notag
\end{align}
Since $T=2R/W_0(Z)$,
\begin{align}
 T\sim\frac{2R}{L_1}\biggl[{}&1+\frac{L_2}{L_1}
 +\frac{L_2(L_2-1)}{L_1^2}
 +\frac{L_2(2L_2^2-5L_2+2)}{2L_1^3}
 \notag\\
 &+\frac{L_2(6L_2^3-26L_2^2+27L_2-6)}{6L_1^4}
 +O\!\left(\frac{L_2^5}{L_1^5}\right)\biggr].
 \label{eq:T-outer-first}
\end{align}
\Cref{eq:um-general,eq:tau-partition} give every omitted coefficient.

\subsection{All coefficients after expanding the inverse blocks}

The compact complete branchwise transseries can be written directly in $R$
and $w=W_0(2R/\e)$ as
\begin{equation}
 \boxed{
 \D_r^{-1}(y)\sim
 \frac{2R_r}{w_r}-1+
 \sum_{n\ge1}\sum_{j=1}^{2n-1}
 c_{n,j}
 \left(\frac{w_r}{2R_r}\right)^{2n-1}
 (w_r+1)^{-j}.}
 \label{eq:compact-full-transseries}
\end{equation}
Together, \cref{eq:closed-c-nj,eq:W-Stirling-expansion} are already a general
formula for every coefficient of the fully elementary expansion.  For
completeness we spell out the final composition step.

Put
\begin{equation}
 v_1=1+u_1=1-L_2,
 \qquad v_m=u_m\quad(m\ge2),
 \label{eq:v-def}
\end{equation}
so that
\begin{equation}
 w+1=L_1\left(1+\sum_{m\ge1}\frac{v_m}{L_1^m}\right).
 \label{eq:wplus1-v}
\end{equation}
For positive integers $p,j$, define
\begin{align}
 (1+\textstyle\sum_{m\ge1}u_mz^m)^p
 &=\sum_{m\ge0}\sigma_m^{(p)}z^m,
 \label{eq:sigma-def}\\
 (1+\textstyle\sum_{m\ge1}v_mz^m)^{-j}
 &=\sum_{m\ge0}\rho_m^{(j)}z^m.
 \label{eq:rho-def}
\end{align}
Their finite partition formulae are
\begin{align}
 \sigma_m^{(p)}={}&
 \sum_{\substack{k_1,\ldots,k_m\ge0\\\sum hk_h=m}}
 \binom{p}{K}\frac{K!}{\prod_hk_h!}
 \prod_{h=1}^m u_h^{k_h},
 \label{eq:sigma-partition}\\
 \rho_m^{(j)}={}&
 \sum_{\substack{k_1,\ldots,k_m\ge0\\\sum hk_h=m}}
 (-1)^K\binom{j+K-1}{K}\frac{K!}{\prod_hk_h!}
 \prod_{h=1}^m v_h^{k_h},
 \label{eq:rho-partition}
\end{align}
where $K=\sum_hk_h$.  It follows that
\begin{equation}
 w^p(w+1)^{-j}
 \sim L_1^{p-j}\sum_{m\ge0}\frac1{L_1^m}
 \sum_{a+b=m}\sigma_a^{(p)}\rho_b^{(j)}.
 \label{eq:w-power-composition}
\end{equation}
Substituting $p=2n-1$ in \cref{eq:w-power-composition} and then using
\cref{eq:compact-full-transseries} gives a finite combinatorial formula for
any prescribed coefficient at any prescribed triple of levels
$(R^{-1},L_1^{-1},L_2)$.

\subsection{Ordering of the scales}

For every fixed $M$,
\begin{equation}
 \frac{T}{L_1^M}\longrightarrow\infty.
 \label{eq:T-beyond-outer}
\end{equation}
Therefore the fixed centring shift and the first Bernoulli blocks obey
\begin{equation}
 \frac{T}{L_1^M}\gg1\gg\frac1{T\Lambda}
 \gg\frac1{T^3\Lambda}\gg\frac1{T^5\Lambda}\gg\cdots.
 \label{eq:nested-scale-ordering}
\end{equation}
This is the reason for retaining $-1$ explicitly in
\cref{eq:branch-inverse-all-orders}.  It is invisible to every finite
outer Lambert expansion, yet dominates the complete inverse-power grid.

\section{The inverse of the periodic OEIS continuation}
\label{sec:periodic-inverse}

The one-function continuation \cref{eq:OEIS-continuation} is useful, but its
inverse cannot be obtained by merely inserting a parity-averaged constant into
the branchwise formula.  The periodic factor creates a new sector before the
Bernoulli grid.

\subsection{Exact displacement equation}

Let
\begin{equation}
 R_*=\log y-C_*,
 \qquad w=\W_0(2R_*/\e),
 \qquad T=\frac{2R_*}{w},
 \qquad \Lambda=w+1,
 \qquad q=T^{-1},
 \qquad \theta=\pi T.
 \label{eq:periodic-core}
\end{equation}
Write the inverse in centred form as
\begin{equation}
 \DO_+^{-1}(y)=T-1+\Delta.
 \label{eq:periodic-inverse-displacement}
\end{equation}
Using the exact identity in \cref{eq:OEIS-centered-exact}, the displacement
satisfies
\begin{equation}
 \Lambda\Delta+q^{-1}\varphi(q\Delta)
 -2a\cos(\theta+\pi\Delta)
 +2\mathcal C(T+\Delta)=0.
 \label{eq:periodic-exact-equation}
\end{equation}
The second term begins with $q\Delta^2/2$; the Borel correction begins with
$-q/6$.  In contrast, the cosine term is $O(1)$.  Thus
$\Delta=O(\Lambda^{-1})$, not $O(q\Lambda^{-1})$.

\subsection{The complete zero-grid log-periodic sector}

At $q=0$, \cref{eq:periodic-exact-equation} reduces to
\begin{equation}
 \Lambda\Delta_0=2a\cos(\theta+\pi\Delta_0).
 \label{eq:periodic-zero-grid}
\end{equation}
For $\Lambda>2\pi a$, the right-hand side divided by $\Lambda$ is a
contraction on the real line, so \cref{eq:periodic-zero-grid} has a unique real
solution.  Its complete inverse-logarithmic expansion has closed Fourier
coefficients.

\begin{theorem}[Closed log-periodic sector]
\label{thm:periodic-zero-sector}
Let $\lambda=\Lambda^{-1}$.  The solution of
\cref{eq:periodic-zero-grid} has the expansion
\begin{equation}
 \Delta_0(\theta,\lambda)
 =\sum_{m\ge1}p_m(\theta)\lambda^m,
 \label{eq:Delta0-expansion}
\end{equation}
where
\begin{align}
 p_m(\theta)
 &={\frac{(2a)^m}{m!}}
 \left.\frac{\dd^{m-1}}{\dd u^{m-1}}
 \cos^m(\theta+\pi u)\right|_{u=0}
 \label{eq:pm-derivative}\\
 &=\frac{a^m}{m!}
 \sum_{j=0}^{m}\binom mj
 \bigl(\ii\pi(m-2j)\bigr)^{m-1}
 \e^{\ii(m-2j)\theta}.
 \label{eq:pm-Fourier}
\end{align}
The second expression is real despite its complex notation.  The first terms
are
\begin{align}
 \Delta_0={}&\frac{2a\cos\theta}{\Lambda}
 -\frac{4\pi a^2\cos\theta\sin\theta}{\Lambda^2}
 \notag\\
 &+\frac{4\pi^2a^3\cos\theta(2-3\cos^2\theta)}{\Lambda^3}
 +O(\Lambda^{-4}).
 \label{eq:Delta0-first}
\end{align}
\end{theorem}

\begin{proof}
\Cref{eq:periodic-zero-grid} is
$\Delta_0=2a\lambda\cos(\theta+\pi\Delta_0)$.  Lagrange--B\"urmann inversion
in the parameter $\lambda$ gives \cref{eq:pm-derivative}.  Expanding
$\cos^m v=2^{-m}\sum_{j=0}^m\binom mj\e^{\ii(m-2j)v}$ and differentiating
gives \cref{eq:pm-Fourier}.  Direct evaluation for $m=1,2,3$ yields
\cref{eq:Delta0-first}.
\end{proof}

The formula in \cref{eq:pm-Fourier} exhibits the precise harmonic content: at
inverse-logarithmic order $m$, only Fourier modes
$m,m-2,\ldots,-m$ occur.  Thus the periodic sector forms a finitely generated
coefficient algebra at every order, even though it lies outside ordinary
polynomial-logarithmic transseries.

\subsection{Coupling the periodic sector to all inverse powers}

Replace the exact $\mathcal C$ in \cref{eq:periodic-exact-equation} by its
Borel-summable asymptotic series and define the formal function
\begin{align}
 \mathcal F(q,u)={}&\Lambda u+
 \sum_{m\ge2}\frac{(-1)^m}{m(m-1)}u^mq^{m-1}
 -2a\cos(\theta+\pi u)
 \notag\\
 &+2\sum_{k\ge1}\beta_kq^{2k-1}(1+qu)^{-(2k-1)}.
 \label{eq:periodic-formal-F}
\end{align}
Seek
\begin{equation}
 \Delta=\Delta_0+\sum_{n\ge1}\Delta_nq^n,
 \label{eq:periodic-q-expansion}
\end{equation}
where each $\Delta_n$ is a formal series in $\Lambda^{-1}$ with trigonometric
polynomial coefficients.  Put
\begin{equation}
 M=\Lambda+2\pi a\sin(\theta+\pi\Delta_0).
 \label{eq:periodic-M}
\end{equation}
The leading term of $M$ is $\Lambda$, so it is invertible in this coefficient
field.

\begin{theorem}[Triangular periodic-grid recurrence]
\label{thm:periodic-grid}
For $n\ge1$, let
\begin{equation}
 \Delta_{<n}(q)=\Delta_0+\sum_{j=1}^{n-1}\Delta_jq^j.
 \label{eq:Delta-less-n}
\end{equation}
Then the unique formal solution of \cref{eq:periodic-formal-F} is generated by
\begin{equation}
 \boxed{
 \Delta_n=-\frac1M\coeff{q}{n}
 \mathcal F\bigl(q,\Delta_{<n}(q)\bigr).}
 \label{eq:periodic-grid-recurrence}
\end{equation}
The first coefficient is
\begin{equation}
 \Delta_1=\frac{1/6-\Delta_0^2/2}{M}.
 \label{eq:Delta1}
\end{equation}
The second is
\begin{align}
 \Delta_2=-\frac1M\biggl[{}&
 a\pi^2\cos(\theta+\pi\Delta_0)\Delta_1^2
 +\Delta_0\Delta_1
 +\frac{\Delta_0-\Delta_0^3}{6}\biggr].
 \label{eq:Delta2}
\end{align}
\end{theorem}

\begin{proof}
The coefficient of a new $\Delta_nq^n$ in the $q^n$ part of
$\mathcal F(q,\Delta)$ is the derivative of the zero-grid function
$\Lambda u-2a\cos(\theta+\pi u)$ at $u=\Delta_0$, namely $M$.  Every other
contribution involves only lower $\Delta_j$.  Solving the coefficient equation
gives \cref{eq:periodic-grid-recurrence}.

At order $q$, the core kernel contributes $\Delta_0^2/2$, and
$2\beta_1=-1/6$ supplies the Bernoulli contribution, proving
\cref{eq:Delta1}.  At order $q^2$, Taylor expansion of the zero-grid function,
the $u^2q/2-u^3q^2/6$ core kernel, and the expansion of
$2\beta_1q(1+qu)^{-1}$ give \cref{eq:Delta2}.
\end{proof}

\begin{corollary}[Asymptotic meaning of the periodic grid]
\label{cor:periodic-asymptotic-meaning}
Fix nonnegative integers $N$ and $M$.  Truncate
\cref{eq:periodic-q-expansion} after $q^N$, and in each coefficient retain the
inverse-logarithmic expansion through $\Lambda^{-M}$.  If the centred forward
series is retained far enough to determine these coefficients, substitution
into \cref{eq:periodic-exact-equation} leaves a residual that
is smaller in the lexicographic scale
\[
 1,\Lambda^{-1},\Lambda^{-2},\ldots;
 \quad q,q\Lambda^{-1},q\Lambda^{-2},\ldots;
 \quad q^2,\ldots .
\]
Since the exact derivative with respect to $\Delta$ equals
$M+O(q)$ and $M\sim\Lambda$, the residual transfers to an inverse error after
division by $\Lambda$.  The assertion is uniform in the phase
$\theta=\pi T$, because all trigonometric coefficients and their derivatives
are bounded.
\end{corollary}

\begin{proof}
The recurrence cancels the formal residual successively in $q$.  At each fixed
$q$-order, Lagrange inversion cancels it successively in $\Lambda^{-1}$.
The exact Borel remainder estimate \cref{thm:remainder-bounds} controls the
replacement of $\mathcal C$ by its finite Bernoulli expansion.  Finally, the
implicit derivative is bounded below by a positive multiple of $\Lambda$ for
large $T$, uniformly in $\theta$, so the mean-value argument used in
\cref{thm:branch-inverse} applies.
\end{proof}

\begin{summarybox}
The eventual inverse of the OEIS continuation is
\begin{equation}
 \DO_+^{-1}(y)
 \sim T-1+\Delta_0(\pi T,\Lambda^{-1})
 +\sum_{n\ge1}\Delta_n(\pi T,\Lambda^{-1})T^{-n}.
 \label{eq:periodic-full-transseries}
\end{equation}
The coefficient $\Delta_0$ is given explicitly by
\cref{eq:pm-derivative,eq:pm-Fourier}; all finer coefficients are generated by
\cref{eq:periodic-grid-recurrence}.  This is a log-periodic transseries over
the monomial group generated by $\Lambda^{-1}$, $T^{-1}$, and
$\e^{\ii\pi T}$.
\end{summarybox}

\subsection{How the parity branches resum the periodic sector}

At an integer $n$, $s=n+1$ and
\begin{equation}
 C_*-a\cos(\pi s)=
 \begin{cases}
 C_0,&n\ \hbox{even},\\
 C_1,&n\ \hbox{odd}.
 \end{cases}
 \label{eq:periodic-freeze}
\end{equation}
The branchwise normalization absorbs this entire constant before the Lambert
inversion.  Hence the large $O(\Lambda^{-1})$ periodic displacement disappears,
and the first remaining correction is the much smaller
$1/(6T\Lambda)$.  In this precise sense the two parity-resolved transseries are
subsequence resummations of the periodic inverse: they build the parity phase
into the core instead of expanding it perturbatively.

\section{Divergence and beyond-all-orders inverse accuracy}
\label{sec:inverse-large-order}

The branchwise inverse series is Poincar\'e asymptotic, not convergent.  An
exact large-order statement follows immediately from the shallow edge of the
coefficient triangle:
\begin{equation}
 [\Lambda^{-1}]\,d_{2n-1}(\Lambda)
 =(-1)^{n+1}\frac{2(2n-2)!\,\etaD(2n)}{\pi^{2n}}
 =(-1)^{n+1}\frac{2(2n-2)!}{\pi^{2n}}
 \left(1+O(2^{-2n})\right).
 \label{eq:d-large-order-leading}
\end{equation}
At this fixed inverse-logarithmic depth, the term
$[\Lambda^{-1}]d_{2n-1}/T^{2n-1}$ is smallest near
$2n\approx\pi T$.  The deeper powers of $\Lambda^{-1}$ are not uniform in
that diagonal limit, so \cref{eq:d-large-order-leading} alone should not be
used as a proof of optimal truncation for the full block polynomial.

A rigorous envelope follows instead from the exact forward remainder.  Divide
the optimally truncated centred logarithmic residual
\cref{eq:least-term-envelope} by the inverse slope
$\tfrac12\Lambda+O(T^{-2})$.  This gives
\begin{equation}
 \operatorname{error}_{\mathrm{opt}}
 =O\!\left(\frac{\e^{-\pi T}}{\sqrt T\,\Lambda}\right).
 \label{eq:inverse-optimal-envelope}
\end{equation}
Thus the exponential envelope is established without assuming any unproved
uniform large-$n$ estimate for the complete polynomials $d_{2n-1}$.

\begin{proposition}[Transport of a forward beyond-all-orders sector]
\label{prop:sector-transport}
Suppose a chosen exponentially improved representation of the centred forward
phase retains the algebraic terms $\beta_1,\ldots,\beta_K$ and has an
additional residual $\mathcal E(s)$.  Let $s_{\mathrm{alg}}$ solve the
correspondingly truncated algebraic inverse problem.  Then the leading induced
inverse correction is
\begin{equation}
 \Delta_{\mathcal E}
 =-\frac{2\mathcal E(s_{\mathrm{alg}})}
 {\log s_{\mathrm{alg}}-2\sum_{k=1}^{K}(2k-1)\beta_k
  s_{\mathrm{alg}}^{-2k}}
 +O\!\left(
 \frac{\sup|F_{\mathrm{alg}}''|\,|\mathcal E|^2}{\Lambda^3}
 +\frac{|\mathcal E\,\mathcal E'|}{\Lambda^2}
 \right),
 \label{eq:sector-transport}
\end{equation}
where $F_{\mathrm{alg}}$ is the doubled algebraic phase and the supremum is
taken between the algebraic and exact inverse points.  The estimate assumes
the displayed denominator is comparable to $\Lambda$ and
$|\mathcal E'|=o(\Lambda)$.
\end{proposition}

\begin{proof}
Let $F(s)$ denote the exact doubled forward phase minus $2R$, and let
$F_{\mathrm{alg}}$ denote the same expression with $\mathcal E$ omitted.
By construction, $F_{\mathrm{alg}}(s_{\mathrm{alg}})=0$ and
$F(s_{\mathrm{alg}})=2\mathcal E(s_{\mathrm{alg}})$.  Taylor expansion
$0=F(s_{\mathrm{alg}}+\Delta)=F(s_{\mathrm{alg}})+F'(s_{\mathrm{alg}})\Delta
+O(F''\Delta^2)$ gives the first Newton correction.  Replacing
$F'(s_{\mathrm{alg}})=F_{\mathrm{alg}}'(s_{\mathrm{alg}})
+2\mathcal E'(s_{\mathrm{alg}})$ by $F_{\mathrm{alg}}'$ and iterating once
produces the two displayed remainder terms and proves
\cref{eq:sector-transport}.
\end{proof}

This proposition is deliberately stated in terms of a supplied forward sector.
The Borel poles at $\pm\ii\pi m$ generate terminant functions whose form and
Stokes multipliers depend on the sector of analytic continuation.  The transport law
in \cref{eq:sector-transport} carries any such rigorously normalized forward
completion without inventing a spurious exponential term.

\section{Complex inverse sheets}
\label{sec:complex}

Although the principal application is the large positive real inverse, the
Lambert normalization also labels complex asymptotic sheets.  Fix a branch of
$\Log y$ and an integer logarithmic winding number $m$.  Put
\begin{equation}
 R_{r,m}=\Log y+2\pi\ii m-C_r.
 \label{eq:complex-R}
\end{equation}
For a Lambert branch $k\in\mathbb Z$, define
\begin{equation}
 w_{r,m,k}=\W_k(2R_{r,m}/\e),
 \qquad
 T_{r,m,k}=\exp(w_{r,m,k}+1)=\frac{2R_{r,m}}{w_{r,m,k}},
 \qquad
 \Lambda_{r,m,k}=w_{r,m,k}+1.
 \label{eq:complex-core}
\end{equation}
Whenever $T_{r,m,k}$ lies in a sector in which the gamma Stirling expansion is
uniform and avoids the poles of the forward branch, the same coefficient
polynomials give the local formal sheet
\begin{equation}
 x_{r,m,k}(y)
 \sim T_{r,m,k}-1+
 \sum_{n\ge1}\frac{d_{2n-1}(\Lambda_{r,m,k})}
 {T_{r,m,k}^{2n-1}}.
 \label{eq:complex-inverse-sheet}
\end{equation}
The pair $(m,k)$ is an asymptotic label, not a global classification: distinct
labels can meet after analytic continuation around gamma poles or logarithmic
cuts.  The expansion in \cref{eq:complex-inverse-sheet} nevertheless gives a convenient
local atlas at large modulus.

\section{Numerical evaluation and certification}
\label{sec:algorithm}

\subsection{Stable evaluation from a logarithmic target}

Directly forming $y=n!!$ quickly overflows fixed-precision arithmetic.  The
inverse should therefore accept $\ell=\log y$ as its primary input.

\begin{algorithm}[Parity-resolved inverse double factorial]
\label{alg:inverse-double-factorial}
Given $\ell=\log y$, parity $r\in\{0,1\}$, and a desired number $N$ of
algebraic correction blocks:
\begin{lstlisting}
C := (log(pi))/2              if r = 0
     (log(2))/2               if r = 1
R := ell - C
w := LambertW(0, 2*R/e)
T := 2*R/w
Lambda := w + 1
x := T - 1
for n := 1,...,N:
    x := x + d[2*n-1](Lambda) / T^(2*n-1)
return x
\end{lstlisting}
The polynomials $d_{2n-1}$ may be precomputed by
\cref{eq:inverse-triangular-recurrence}, or individual scalar coefficients may
be generated by \cref{eq:closed-c-nj}.
\end{algorithm}

For full working precision, use the transseries approximation as a starting
point for Newton's method applied to the logarithmic residual
\begin{equation}
 F_r(x)=\log\kappa_r+\frac{x}{2}\log2
 +\log\Gamma\!\left(\frac{x}{2}+1\right)-\ell.
 \label{eq:Newton-residual}
\end{equation}
Its derivative is the positive quantity in \cref{eq:Dr-logder-exact}, so the
Newton update is
\begin{equation}
 x_{\mathrm{new}}=x-
 \frac{F_r(x)}{\tfrac12\left(\log2+\psi(x/2+1)\right)}.
 \label{eq:Newton-update}
\end{equation}
A few asymptotic blocks normally place the initial value well inside the
quadratic convergence region.

For the OEIS interpolation, add
$\tfrac14(1-\cos\pi x)\log(2/\pi)$ to $F_0(x)$ and its derivative to the
Newton denominator.  The transseries in \cref{eq:periodic-full-transseries} gives the
appropriate initial approximation.

\subsection{A posteriori inverse certificate}

Let $\widetilde x$ be any approximation and let
$\rho=F_r(\widetilde x)$.  If an interval $I$ containing both
$\widetilde x$ and the exact inverse has
\begin{equation}
 m_I=\inf_{x\in I}F_r'(x)>0,
 \label{eq:slope-lower-bound}
\end{equation}
then the mean-value theorem gives
\begin{equation}
 |\D_r^{-1}(y)-\widetilde x|\le\frac{|\rho|}{m_I}.
 \label{eq:a-posteriori-bound}
\end{equation}
Because $F_r'$ is increasing, a convenient lower bound is its value at the left
endpoint of $I$.  Combining \cref{eq:a-posteriori-bound} with the parity floor
in \cref{eq:discrete-parity-inverse} certifies the exact integer generalized
inverse whenever the resulting interval does not cross a same-parity integer.

\subsection{Numerical validation}

To test the expansion without root-finding ambiguity, choose an exact index
$n$, set $\ell=\log\D_r(n)$, reconstruct the Lambert variables from $\ell$,
and compare each truncation with $n$.  The following values were computed with
80-decimal arithmetic.  Each entry is approximation minus exact index; $E_0$
uses only $T-1$, and $E_j$ successively includes $d_1,d_3,\ldots,d_{2j-1}$.

\begin{table}[htbp]
\centering
\caption{Errors for the even branch $r=0$.}
\label{tab:even-errors}
\scriptsize
\resizebox{\textwidth}{!}{%
\begin{tabular}{r@{\quad}rrrrr}
\toprule
$n$ & $E_0$ & $E_1$ & $E_2$ & $E_3$ & $E_4$\\
\midrule
10  & $-6.3073664\times10^{-3}$ & $ 1.6443978\times10^{-5}$
    & $-1.6096633\times10^{-7}$ & $ 3.6889599\times10^{-9}$
    & $-1.6035926\times10^{-10}$\\
20  & $-2.6054931\times10^{-3}$ & $ 1.7519735\times10^{-6}$
    & $-4.7751372\times10^{-9}$ & $ 3.1028710\times10^{-11}$
    & $-3.8172037\times10^{-13}$\\
50  & $-8.3108660\times10^{-4}$ & $ 8.9793611\times10^{-8}$
    & $-4.1940787\times10^{-11}$ & $ 4.7155771\times10^{-14}$
    & $-9.9985171\times10^{-17}$\\
100 & $-3.5754817\times10^{-4}$ & $ 9.5804953\times10^{-9}$
    & $-1.1469146\times10^{-12}$ & $ 3.3159325\times10^{-16}$
    & $-1.8032908\times10^{-19}$\\
\bottomrule
\end{tabular}%
}
\end{table}

\begin{table}[htbp]
\centering
\caption{Errors for the odd branch $r=1$.}
\label{tab:odd-errors}
\scriptsize
\resizebox{\textwidth}{!}{%
\begin{tabular}{r@{\quad}rrrrr}
\toprule
$n$ & $E_0$ & $E_1$ & $E_2$ & $E_3$ & $E_4$\\
\midrule
9  & $-7.2226927\times10^{-3}$ & $ 2.3054068\times10^{-5}$
   & $-2.7222582\times10^{-7}$ & $ 7.4928917\times10^{-9}$
   & $-3.9120795\times10^{-10}$\\
19 & $-2.7801819\times10^{-3}$ & $ 2.0689304\times10^{-6}$
   & $-6.2119116\times10^{-9}$ & $ 4.4426296\times10^{-11}$
   & $-6.0165406\times10^{-13}$\\
49 & $-8.5199643\times10^{-4}$ & $ 9.5862235\times10^{-8}$
   & $-4.6575971\times10^{-11}$ & $ 5.4466224\times10^{-14}$
   & $-1.2012435\times10^{-16}$\\
99 & $-3.6190377\times10^{-4}$ & $ 9.8955366\times10^{-9}$
   & $-1.2083617\times10^{-12}$ & $ 3.5634509\times10^{-16}$
   & $-1.9767173\times10^{-19}$\\
\bottomrule
\end{tabular}%
}
\end{table}

The parity dependence is confined to the Lambert core.  Once the proper
constant $C_r$ is used, the same sequence of coefficient polynomials produces
nearly identical convergence patterns on the two subsequences.

\section{Checks, alternative derivations, and implementation notes}
\label{sec:checks}

\subsection{Formal residual check}

A generated block list should always be checked by substitution rather than by
comparing displayed numerators.  If
\begin{equation}
 E_N(q)=\sum_{n=1}^{N}d_{2n-1}(\Lambda)q^{2n},
 \label{eq:EN-check}
\end{equation}
then the exact formal certificate is
\begin{equation}
 \Lambda E_N+\varphi(E_N)
 +\sum_{k=1}^{N+1}\alpha_kq^{2k}(1+E_N)^{-(2k-1)}
 =O(q^{2N+2}).
 \label{eq:formal-residual-certificate}
\end{equation}
The coefficient of $q^{2N+2}$ predicts the next block after division by
$-\Lambda$.  This identity catches sign, indexing, and centring mistakes at
once.

\subsection{Independent checks}

The formulae in this article admit four independent comparisons.
\begin{enumerate}
\item Expanding the exact gamma ratio \cref{eq:Dr-ratio} reproduces
      $\beta_k$ directly.
\item Expanding the Borel kernels \cref{eq:U-partial-fractions,eq:C-partial-fractions}
      reproduces $\delta_k$ and $\beta_k$ from zeta and eta values.
\item The triangular inverse recurrence and the closed
      Lagrange--Bell formula agree coefficient by coefficient.
\item The scaling identity in \cref{eq:inverse-gamma-scaling} compares every
      inverse block with the independently derived half-shifted inverse-Gamma
      expansion in \cite{ProveItGammaInverse}.
\end{enumerate}
The displayed coefficients through $d_{11}$ pass all four checks.

\subsection{Truncation discipline}

At recurrence step $n$, every formal operation may be truncated above
$q^{2n}$.  In particular, there is no reason to form a full symbolic logarithm
or a full negative binomial series.  A practical implementation maintains a
vector of coefficients in $q^2$ whose entries are rational functions of
$\Lambda$; multiplication is truncated convolution.  The formal residual
\cref{eq:formal-residual-certificate} should be evaluated after each step.

For high $n$, it is often more efficient to represent a rational function in
$z=\Lambda^{-1}$ as a dense polynomial and postpone common-denominator
normalization.  The degree bound $2n-1$ gives an exact allocation size.  The
\Cref{thm:sign-law} provides a cheap additional invariant: after
multiplication by $(-1)^{n+1}$, every coefficient must be positive.

\section{Conclusions}
\label{sec:conclusion}

The double factorial has a particularly transparent all-orders asymptotic
structure once the two correct normalizations are made.

First, parity must be resolved.  The branches $\D_0$ and $\D_1$ share the same
gamma factor and differ only by a target constant.  This isolates parity in
the Lambert core and leaves a universal correction calculus.  For the integer
sequence, the equivalent leading amplitude is the periodic transmonomial
$A_n=(2\pi)^{1/4}(\pi/2)^{(-1)^n/4}$.

Second, the shift $s=x+1$ is the natural inversion coordinate.  It converts the
forward logarithm to
\[
 \frac{s}{2}(\log s-1)+C_r+\text{odd inverse powers},
\]
so the exact dominant inverse is $2R/W_0(2R/\e)$ and the correction support is
an even grid in the relative displacement.  Every inverse block is a finite
polynomial in $1/\log T$.

The coefficient problem is completely solved in two complementary ways.  The
triangular recurrence gives an efficient constructive algorithm, while the
Lagrange--B\"urmann formula in \cref{eq:closed-c-nj} expresses every scalar
coefficient as a finite sum of partial Bell polynomials.  It yields degree and
support bounds, boundary coefficients, and the coefficientwise alternating
sign law.  Unsigned Stirling numbers then generate the outer expansion of the
Lambert core, so the complete elementary transseries has explicit
combinatorial coefficients at every level.

The exact Borel kernels
\[
 \frac{\coth t}{2t}-\frac1{2t^2},
 \qquad
 \frac1{2t\sinh t}-\frac1{2t^2}
\]
complete the direct asymptotic series analytically.  They give rigorous
first-omitted-term bounds on the positive axis and identify the pole lattice
responsible for divergence and hyperasymptotic corrections.  The inverse
transport formula then converts any normalized forward exponential sector
into its inverse counterpart.

Finally, the single OEIS interpolation demonstrates why the interpolation
choice cannot be ignored.  Its periodic factor survives at leading logarithmic
order and produces a bounded but rapidly oscillating $1/\log T$ sector.  The
closed Fourier--Lagrange formula and the periodic-grid recurrence provide its
complete inverse transseries.  The parity-resolved expansions are not competing
answers: they are the cleaner subsequence-resummed forms appropriate when the
original integer parity is known.

\appendix

\section{Coefficient data in separated inverse-logarithmic form}
\label{app:separated}

For some comparisons it is useful to separate each block into powers of
$\Lambda^{-1}$.  The first four are
\begin{align}
 d_1={}&\frac1{6\Lambda},
 \notag\\
 d_3={}&-\frac7{180\Lambda}
         -\frac1{36\Lambda^2}
         -\frac1{72\Lambda^3},
 \label{eq:d3-separated}\\
 d_5={}&\frac{31}{630\Lambda}
 +\frac7{270\Lambda^2}
 +\frac{17}{1080\Lambda^3}
 +\frac5{648\Lambda^4}
 +\frac1{432\Lambda^5},
 \label{eq:d5-separated}\\
 d_7={}&-\frac{127}{840\Lambda}
 -\frac{4063}{75600\Lambda^2}
 -\frac{11413}{453600\Lambda^3}
 -\frac1{80\Lambda^4}
 \notag\\
 &-\frac{157}{25920\Lambda^5}
 -\frac{35}{15552\Lambda^6}
 -\frac5{10368\Lambda^7}.
 \label{eq:d7-separated}
\end{align}
The compact common-denominator forms in \cref{eq:d3,eq:d5,eq:d7} are less
susceptible to transcription error; the separated form makes
\cref{eq:d-leading-slope} visible.

\section{A generic coefficient-generation pseudocode}
\label{app:pseudocode}

The following pseudocode implements \cref{eq:inverse-triangular-recurrence}
on the variable $z=q^2$.
\begin{lstlisting}
input: N, symbolic Lambda
E := 0
for n := 1,...,N:
    H := (1+E)*log(1+E) - E
    for k := 1,...,n:
        H := H + alpha[k]*z^k*(1+E)^(-(2*k-1))
    H := truncate(H, z^(n+1))
    e[n] := -coefficient(H, z^n)/Lambda
    E := E + e[n]*z^n
    assert truncate(residual(E), z^(n+1)) = 0
return e[1],...,e[N]
\end{lstlisting}
All logarithms and powers here are formal truncated series.  The assertion uses
\cref{eq:formal-residual-certificate}.

\section{Notation crosswalk}
\label{app:crosswalk}

\begin{center}
\begin{tabularx}{\textwidth}{@{}lX@{}}
\toprule
Symbol & Meaning\\
\midrule
$\D_r(x)$ & parity-resolved gamma interpolation, $r=0$ even and $r=1$ odd\\
$\DO(x)$ & periodic analytic continuation agreeing with all integer values\\
$s=x+1$ & optimal centred forward variable\\
$C_0,C_1,C_*$ & even, odd, and parity-mean logarithmic constants\\
$\delta_k$ & uncentred logarithmic Stirling coefficients\\
$\beta_k$ & centred logarithmic coefficients\\
$\alpha_k=2\beta_k$ & coefficients in the doubled inverse phase\\
$R_r$ & $\log y-C_r$\\
$w_r$ & $W_0(2R_r/\e)$\\
$T_r$ & exact Lambert core $2R_r/w_r$\\
$\Lambda_r$ & core slope $w_r+1=\log T_r$\\
$d_{2n-1}$ & branchwise inverse correction polynomial\\
$c_{n,j}$ & coefficient of $\Lambda^{-j}$ in $d_{2n-1}$\\
$\Delta_0,\Delta_n$ & periodic inverse sectors for $\DO_+^{-1}$\\
\bottomrule
\end{tabularx}
\end{center}

\begin{thebibliography}{99}

\bibitem{OEIS}
OEIS Foundation Inc.,
\emph{The On-Line Encyclopedia of Integer Sequences},
entry A006882, ``Double factorials $n!!$,''
\url{https://oeis.org/A006882}.

\bibitem{DLMFGamma}
NIST Digital Library of Mathematical Functions,
\S5.11, ``Asymptotic Expansions'' in Chapter 5, Gamma Function,
\url{https://dlmf.nist.gov/5.11}.

\bibitem{DLMFDuplication}
NIST Digital Library of Mathematical Functions,
\S5.5, ``Functional Relations'' in Chapter 5, Gamma Function,
\url{https://dlmf.nist.gov/5.5}.

\bibitem{DLMFW}
NIST Digital Library of Mathematical Functions,
\S4.13, ``Lambert $W$-Function,''
\url{https://dlmf.nist.gov/4.13}.

\bibitem{Corless}
R.~M. Corless, G.~H. Gonnet, D.~E.~G. Hare, D.~J. Jeffrey, and D.~E. Knuth,
On the Lambert $W$ function,
\emph{Advances in Computational Mathematics} \textbf{5} (1996), 329--359.
\url{https://doi.org/10.1007/BF02124750}.

\bibitem{Olver}
F.~W.~J. Olver,
\emph{Asymptotics and Special Functions},
AKP Classics, A K Peters, Wellesley, MA, 1997.

\bibitem{Comtet}
L. Comtet,
\emph{Advanced Combinatorics: The Art of Finite and Infinite Expansions},
Reidel, Dordrecht, 1974.

\bibitem{FlajoletSedgewick}
P. Flajolet and R. Sedgewick,
\emph{Analytic Combinatorics},
Cambridge University Press, Cambridge, 2009.

\bibitem{Meserve}
B.~E. Meserve,
Double factorials,
\emph{American Mathematical Monthly} \textbf{55} (1948), 425--426.

\bibitem{ParisKaminski}
R.~B. Paris and D. Kaminski,
\emph{Asymptotics and Mellin--Barnes Integrals},
Cambridge University Press, Cambridge, 2001.

\bibitem{ProveItCCC}
V. Reshetnikov,
\emph{Combinatorial Coefficient Calculus},
ProveIt semi-formalized research-frontier draft, September 2026,
repository revision \texttt{b150e77f2a03d776201159d5b45f44c50c13c268}.
\url{https://github.com/VladimirReshetnikov/ProveIt/tree/main/Analysis/FabiusFunction/docs/semi-formalized-research-frontiers/drafts/combinatorial-coefficient-calculus/Combinatorial_Coefficient_Calculus}.

\bibitem{ProveItGammaInverse}
V. Reshetnikov,
\emph{Asymptotic Inversion of the Gamma and Barnes $G$-Functions:
Lambert-Core Transseries and a General Reversion Calculus},
ProveIt semi-formalized research-frontier draft, 3 September 2026,
repository revision \texttt{b150e77f2a03d776201159d5b45f44c50c13c268}.
\url{https://github.com/VladimirReshetnikov/ProveIt/tree/main/Analysis/FabiusFunction/docs/semi-formalized-research-frontiers/drafts/series-and-transseries/special-function-inversion/Asymptotic_Inversion_Gamma_Barnes_G}.

\end{thebibliography}

\end{document}
